\documentclass{article}
\usepackage{graphicx}
\usepackage{amsmath}
\usepackage{amssymb}
\usepackage{booktabs}
\usepackage[colorlinks=true,
            linkcolor=blue,
            citecolor=blue,
            urlcolor=blue]{hyperref}
\usepackage{url}
\usepackage{algorithm}
\usepackage{algpseudocode}
\usepackage{float}
\usepackage{tikz}
\usetikzlibrary{arrows.meta, positioning}
\numberwithin{equation}{subsection}
% Mathematical theorem environments
\usepackage{amsthm}

\newtheorem{theorem}{Theorem}[section]
\newtheorem{lemma}[theorem]{Lemma}
\newtheorem{proposition}[theorem]{Proposition}
\newtheorem{corollary}[theorem]{Corollary}

\theoremstyle{definition}
\newtheorem{definition}[theorem]{Definition}
\newtheorem{assumption}[theorem]{Assumption}

\theoremstyle{remark}
\newtheorem{remark}[theorem]{Remark}

\title{Homogeneity-Based Coordinates and Adaptive Grids for Tensor-Train Surrogates in Multi-Asset Option Pricing}
\author{Guillaume Melis^1 \\Antonin Decouvelaere^1 \\Ichak Koutchoukali^2}
\date{August 2026}

\begin{document}

\maketitle
\textit{Keywords: Tensor Network, adaptive discretization, option pricing}
\begin{abstract}
\textit{This paper develops a coordinate-adaptive discretization for the Singularity
Tensor Network pricing framework of Gribben et al.\ \cite{Gribben2026STNGPR}. A tensor-train surrogate reconstructs a discretized price surface from a small number of black-box pricer calls and interpolates between nodes.
The original uniform grid uses absolute strike as a tensor coordinate and may
therefore resolve the high-curvature at-the-money region unevenly across spot
configurations. We exploit the positive homogeneity of basket-option prices to
factor out the basket scale and replace absolute strike with basket
log-moneyness. A hyperbolic-sine grid then concentrates interpolation nodes near
the at-the-money region, while an additional transformation refines short
maturities, without changing the underlying Quantized Tensor-Train
architecture. In the European benchmark, the resulting representation reduces
the effective rank from \(21.45\) to \(7.35\), the price mean absolute error from
\(0.648\) to \(0.0831\), and the measured prediction time by \(83.3\%\), reaching
the interpolation-error floor of the adaptive grid. For American options, the
same coordinate design lowers the mean error and its dispersion across
reconstruction seeds, although the effective ranks remain broadly unchanged
and accuracy is partly limited by the bias and sampling noise of the least-squares
Monte Carlo pricing oracle. Once built, the resulting grid supports near-instantaneous repricing across strikes and maturities via interpolation, making it well suited for fast, on-demand pricing in trading workflows. These results show that coordinate design is a
structural component of accurate and efficient Tensor-Train pricing surrogates.}
\end{abstract}

\newpage
\section{Introduction}
%\subsection{Motivation}
Derivative pricing becomes computationally expensive when closed-form solutions are unavailable, particularly for path-dependent, early-exercise, or multi-asset payoffs.Basket options are a central example: beyond three or four underlyings, standard grid-based PDE solvers and lattice methods scale exponentially in the number of assets, and even Monte Carlo methods become expensive to run repeatedly. While numerical methods remain tractable for a single valuation, their repeated use for large portfolios, Greeks, calibration, or market-risk calculations may require millions of costly valuations. Surrogate models can help reduce this burden by learning a discretized price surface from a small number of black-box pricer calls and interpolates prices between grid nodes, avoiding both the need to solve the pricing problem directly at every query and the exponential grid size a naive discretization would require.
However, conventional surrogates approaches remain affected by the curse of dimensionality.\\
\\
Gribben et al.\ \cite{Gribben2026STNGPR} introduced STN-GPR, a framework combining Gaussian process regression with Tensor-Train representations constructed through TT-cross approximation \cite{oseledets2010ttcross,oseledets2011tensorTrain}. Each discretized coordinate is further split into binary modes, so that the pricing tensor is represented in Quantized Tensor-Train (QTT) form, in order to reduce storage complexity. However, the framework's efficiency also depends on how that discretization is chosen. Uniform grids may oversample smooth regions while providing insufficient resolution near at-the-money configurations, short maturities, or American exercise frontiers, where option prices exhibit stronger curvature. This motivates the use of homogeneous log-moneyness coordinates and convexity-guided non-uniform QTT grids to improve approximation accuracy without increasing the overall grid size.\\
%\subsection{Limitations of uniform high-dimensional pricing grids}
\\Tensor-Train and TT-cross methods reduce storage and the number of pricer
calls by exploiting low-rank structure, but they inherit the discretization
error of the underlying grid: with $n_j$ nodes along dimension $j$, a full
tensor-product grid contains
\[
N_{\mathrm{full}}=\prod_{j=1}^{d} n_j
\]
points, and refining it uniformly removes that error only by enlarging the
tensorized domain and raising the cost of TT-cross construction. We therefore
hold the QTT-compatible node count fixed and act on node placement instead. This allows the grid to concentrate resolution in financially relevant regions without increasing its nominal
tensor size.\\
\\
% A uniform grid implicitly assumes that the pricing function requires the same resolution throughout the parameter domain. This assumption is poorly suited to option-pricing surfaces. Prices are generally smooth and nearly linear far from the strike, whereas stronger curvature appears near at-the-money configurations. Similar local variations arise for short maturities and around the early-exercise frontier of American options. Consequently, a uniform grid allocates many nodes to smooth regions while providing insufficient resolution where the price and its derivatives change most rapidly.\\
% \\
% Increasing the number of uniform nodes can reduce this error, but it also enlarges the tensorized domain and raises the cost of TT-cross construction. A more efficient approach is therefore to preserve the QTT-compatible number of nodes while modifying their physical locations. This allows the grid to concentrate resolution in financially relevant regions without increasing its nominal tensor size.\\
% \\
%\subsection{Contributions}
The contributions of this work are fourfold:
\begin{enumerate}
\item We reproduce the STN-GPR framework for European and American multi-asset options and evaluate it against Gaussian process regression under comparable training budgets.
\item We exploit the homogeneity of option prices to express the pricing surface in normalized log-moneyness coordinates, reducing scale dependence and concentrating the approximation on financially meaningful variables.
\item We introduce a convexity-guided non-uniform grid that concentrates nodes around the at-the-money region while preserving the power-of-two structure required for QTT tensorization. Its performance is assessed through fixed-budget ablation and robustness experiments against the corresponding uniform grid.
\item We provide a detailed error analysis using analytical European benchmarks and independently generated American LSMC labels. For American options, we additionally separate surrogate approximation error from Monte Carlo label noise and examine the accuracy of the early-exercise region.
\end{enumerate}
%\subsection{Outline of the paper}
The remainder of the paper is organized as follows. Section~2 reviews the STN-GPR framework and identifies its main limitations. Section~3 introduces the homogeneity-based transformation and log-moneyness coordinates, while Section~4 presents the convexity-guided adaptive QTT grid. Section~5 describes the error decomposition and validation protocol. Sections~6 report the European and American option experiments, respectively, including the effects of LSMC noise on the exercise frontier. Finally, Section~7 discusses the results and concludes the paper.
\section{The STN-GPR Framework and Its Limitations}
\label{sec:stn-framework}
\subsection{Tensor-Train and Quantized Tensor-Train Representations}
\label{subsec:tt-qtt}
Let \(d\in\mathbb{N}^{*}\) denote the number of pricing variables and let
\(V:\Omega\subset\mathbb{R}^{d}\to\mathbb{R}\) be the associated pricing function.
For each coordinate \(k\in\{1,\ldots,d\}\), consider a grid of
\(n_k\in\mathbb{N}\), \(n_k\geq 2\), nodes denoted by
\(\{x_{k,i_k}\}_{i_k=0}^{n_k-1}\).
The discretization of \(V\) on the resulting Cartesian grid defines the tensor
\begin{equation}
\mathcal{A}_{i_1,\ldots,i_d}
=
V\!\left(
x_{1,i_1},\ldots,x_{d,i_d}
\right)
\end{equation}
where
\begin{equation}
\mathcal{A}\in\mathbb{R}^{n_1\times\cdots\times n_d},
\qquad
0\leq i_k<n_k,
\qquad
k=1,\ldots,d
\end{equation}
A tensor is formally an element of a tensor-product space; once bases are fixed,
it can be represented as a multidimensional array. A tensor network represents
such an array as a collection of lower-order tensors whose shared indices are
contracted. The Tensor-Train (TT) representation is the particular tensor-network
topology in which the component tensors form a one-dimensional chain
\cite{oseledets2011tensorTrain}. It is equivalent to a matrix product state with
open boundary conditions \cite{orus2014tensorNetworks}.\\
\\
%\paragraph{Tensor-Train representation.}
A \textbf{Tensor-Train (TT)} representation of
\(\mathcal{A}\in\mathbb{R}^{n_1\times\cdots\times n_d}\)
expresses its entries as
\begin{equation}
\mathcal{A}_{i_1,\ldots,i_d}
=
\sum_{\alpha_1=1}^{r_1}
\cdots
\sum_{\alpha_{d-1}=1}^{r_{d-1}}
G^{(1)}_{1,i_1,\alpha_1}
G^{(2)}_{\alpha_1,i_2,\alpha_2}
\cdots
G^{(d)}_{\alpha_{d-1},i_d,1}
\end{equation}
where the TT cores satisfy
\begin{equation}
G^{(k)}
\in
\mathbb{R}^{r_{k-1}\times n_k\times r_k},
\qquad
r_0=r_d=1
\end{equation}
The integers
\(\boldsymbol{r}=(r_0,r_1,\ldots,r_d)\)
are the TT ranks. Their minimal values are determined by the ranks of the
successive tensor unfoldings
\begin{equation}
r_k
=
\operatorname{rank}
\left(
\mathcal{A}^{\langle k\rangle}
\right),
\qquad
k=1,\ldots,d-1
\label{eq:tt-unfolding-rank}
\end{equation}
where
\begin{equation}
\mathcal{A}^{\langle k\rangle}
\in
\mathbb{R}^{(n_1\cdots n_k)\times(n_{k+1}\cdots n_d)}
\end{equation}
denotes the matricization that separates the first \(k\) modes from the
remaining \(d-k\) modes.\\
\\
Every finite-dimensional tensor admits an exact TT representation, although
the representation is computationally advantageous only when the TT ranks
remain sufficiently small. Its storage cost is
\begin{equation}
\sum_{k=1}^{d} n_k r_{k-1}r_k
\label{eq:tt-storage}
\end{equation}
Defining
\begin{equation}
n=\max_{1\leq k\leq d}n_k,
\qquad
r=\max_{0\leq k\leq d}r_k
\end{equation}
gives the upper bound
\begin{equation}
\sum_{k=1}^{d} n_k r_{k-1}r_k
\leq dnr^2
\end{equation}
Thus, TT storage scales as
\(\mathcal{O}(dnr^2)\), compared with
\(\mathcal{O}(n^d)\) for a full tensor when \(n_k=n\) for all \(k\).
%In practice, an exact representation is unnecessary. Given a tolerance
%$\varepsilon>0$, the objective is to construct an approximation
%$\widehat{\mathcal{A}}$ satisfying
%\begin{equation}
%\frac{\left\|\mathcal{A}-\widehat{\mathcal{A}} \right\|_F}{\left\| \mathcal{A} \right\|_F} \leq \varepsilon
%\end{equation}
%with ranks as small as possible. The efficiency of the representation therefore
%depends on approximate low-rank structure in the selected TT ordering. This
%property is not implied by smoothness alone: the ranks also depend on parameter
%interactions, coordinate transformations, mode ordering, and localized
%structures such as payoff kinks or exercise boundaries.
%\paragraph{Quantized Tensor-Train representation.}
\\
\\
The \textbf{Quantized Tensor-Train (QTT)} representation further decomposes
each discrete mode into a collection of binary modes. Assume, for simplicity,
that
\begin{equation}
n_k=2^{L_k},
\qquad
L_k=\log_2 n_k
\end{equation}
Each grid index \(i_k\in\{0,\ldots,n_k-1\}\) can then be expressed through its
binary expansion
\begin{equation}
i_k
=
\sum_{\ell=1}^{L_k}
2^{L_k-\ell}b_{k,\ell},
\qquad
b_{k,\ell}\in\{0,1\}
\end{equation}
The original tensor is thereby reshaped into the higher-order binary tensor
defined by
\begin{equation}
\mathcal{A}_{i_1,\ldots,i_d}
=
\widetilde{\mathcal{A}}_{b_{1,1},\ldots,b_{1,L_1},
\ldots,b_{d,1},\ldots,b_{d,L_d}}
\end{equation}
whose order is
\begin{equation}
L
=
\sum_{k=1}^{d}L_k
\end{equation}
A TT decomposition is then applied to
\(\widetilde{\mathcal{A}}\). If its QTT ranks are bounded by
\(r_{\mathrm{QTT}}\), the storage complexity becomes
\begin{equation}
\mathcal{O}\!\left(
r_{\mathrm{QTT}}^2
\sum_{k=1}^{d}\log_2 n_k
\right)
\end{equation}
For \(n_k=n\), this reduces to
\begin{equation}
\mathcal{O}\!\left(
d\,r_{\mathrm{QTT}}^2\log_2 n
\right)
\end{equation}
Quantization is an exact reshaping of the tensor indices and therefore
introduces no approximation by itself. Compression occurs only when the
reshaped tensor admits a low-rank TT representation. Moreover, low TT ranks
for the original tensor do not necessarily imply low QTT ranks after
quantization. The latter depend on the multiscale structure of the indexed
data, the binary encoding convention, and the ordering of the quantized modes.
Throughout this work, the binary modes associated with each pricing variable
are kept contiguous and ordered from the most significant to the least
significant bit.\\
\\
Quantization acts on the discrete indices rather than on the physical
locations of the grid nodes. These locations may therefore be defined through
a strictly increasing mapping
\begin{equation}
x_{k,i}
=
\phi_k\!\left(
\frac{i}{n_k-1}
\right),
\qquad
i=0,\ldots,n_k-1
\label{eq:grid_nodes}
\end{equation}
where \(\phi_k:[0,1]\to\Theta_k\). Affine mappings generate uniform grids,
whereas nonlinear mappings generate non-uniform grids without altering the
binary tensorization. This flexibility will later be used to adapt the
physical grid to the geometry of the pricing function.\\
\\
The TT and QTT representations compress only the values defined at the grid
nodes. They neither determine these values nor define the pricing function
away from the grid. The former problem is addressed through TT-cross
reconstruction, while the latter requires an off-grid interpolation
procedure.
\subsection{TT-Cross Reconstruction}
\label{subsec:tt-cross}
Evaluating every entry of
\(\mathcal{A}\in\mathbb{R}^{n_1\times\cdots\times n_d}\)
requires
\begin{equation}
N_{\mathrm{full}} = \prod_{k=1}^{d} n_k
\end{equation}
calls to the pricing function. TT-cross avoids this exhaustive construction by
accessing individual tensor entries through the oracle
\begin{equation}
(i_1,\ldots,i_d) \longmapsto \mathcal{A}_{i_1,\ldots,i_d} = V(x_{1,i_1},\ldots,x_{d,i_d})
\end{equation}
and adaptively selecting a subset of entries from which a TT approximation is
constructed \cite{oseledets2010ttcross}.\\
\\
TT-cross extends matrix skeleton interpolation to higher-order tensors. To
introduce the underlying principle, let
\(A\in\mathbb{R}^{m\times n}\), where \(m,n\in\mathbb{N}^{*}\), and choose an
integer \(r\leq\min(m,n)\). Let
\begin{equation}
I\subseteq\{1,\ldots,m\},
\qquad
J\subseteq\{1,\ldots,n\},
\qquad
|I|=|J|=r
\end{equation}
If the pivot submatrix \(A(I,J)\in\mathbb{R}^{r\times r}\) is nonsingular, the
associated skeleton approximation is
\begin{equation}
A \approx A(:,J)\,A(I,J)^{-1}A(I,:)
\end{equation}
where \(A(:,J)\in\mathbb{R}^{m\times r}\) contains the selected columns and
\(A(I,:)\in\mathbb{R}^{r\times n}\) contains the selected rows.
If \(\operatorname{rank}(A)=r\), then any nonsingular \(r\times r\) pivot
submatrix yields the exact identity
\begin{equation}
A = A(:,J)\,A(I,J)^{-1}A(I,:)
\end{equation}
When \(\operatorname{rank}(A)>r\), the formula instead defines a rank-\(r\)
approximation whose accuracy and numerical stability depend on the selected
rows and columns. In particular, an ill-conditioned pivot matrix may amplify
the approximation and rounding errors through \(A(I,J)^{-1}\).\\
\\
A common selection principle is therefore to seek a pivot submatrix with large
volume, defined by
\begin{equation}
\operatorname{vol}\!\left(A(I,J)\right) = \left|\det A(I,J)\right|
\end{equation}
A globally maximal-volume submatrix is generally too expensive to compute.
Practical cross algorithms instead use greedy \emph{maxvol} procedures to find
a locally dominant submatrix, providing a tractable and numerically stable
choice of interpolation pivots.\\
\\
The same principle extends recursively to multidimensional tensors. For each
interface $k=1,\ldots,d-1$, TT-cross introduces left multi-indices
\[
\mathcal{I}_{\leq k}
\subset
\{0,\ldots,n_1-1\}
\times\cdots\times
\{0,\ldots,n_k-1\}
\]
and right multi-indices
\[
\mathcal{J}_{>k}
\subset
\{0,\ldots,n_{k+1}-1\}
\times\cdots\times
\{0,\ldots,n_d-1\}
\]
with $\left| \mathcal{I}_{\leq k} \right| = \left| \mathcal{J}_{>k} \right| = r_k$. We also set $\mathcal{I}_{<1}=\{\varnothing\}$ and $\mathcal{J}_{>d}=\{\varnothing\}$. The index sets are nested: each element of $\mathcal{I}_{\leq k}$ extends an element of $\mathcal{I}_{<k}$ by one physical index, with an analogous construction for the right indices. At interface $k$, define the pivot matrix
\begin{equation}
P_k = \mathcal{A} \left( \mathcal{I}_{\leq k}, \mathcal{J}_{>k} \right) \in \mathbb{R}^{r_k\times r_k}
\end{equation}
Assuming nonsingular pivots, the tensor is reconstructed from cross sections
\[
\mathcal{A}
\left(
\mathcal{I}_{<k},
i_k,
\mathcal{J}_{>k}
\right)
\]
and the corresponding TT cores can be written schematically as
\begin{equation}
G^{(k)}(i_k) = \mathcal{A} \left( \mathcal{I}_{<k}, i_k, \mathcal{J}_{>k} \right) P_k^{-1}, \qquad k=1,\ldots,d-1
\end{equation}
with $G^{(d)}(i_d) = \mathcal{A} \left( \mathcal{I}_{<d}, i_d \right)$. The resulting approximation takes the standard TT form
\begin{equation}
\widehat{\mathcal{A}}(i_1,\ldots,i_d) = G^{(1)}(i_1) G^{(2)}(i_2) \cdots G^{(d)}(i_d)
\end{equation}
For an exact tensor of TT ranks $(r_1,\ldots,r_{d-1})$, appropriately selected
nested cross indices recover the tensor exactly in exact arithmetic. In
practice, pivot systems are solved through stable matrix factorizations rather
than by explicitly computing $P_k^{-1}$.
The interpolation sets are constructed iteratively through alternating
left-to-right and right-to-left sweeps over the tensor modes. At mode \(k\),
TT-cross evaluates a subtensor indexed by the current left and right
interpolation sets, reshapes it into a matrix, and applies a rank-revealing
factorization followed by a \emph{maxvol} selection. The selected rows or
columns update the interpolation sets and the corresponding TT core.
Alternating sweeps refine these quantities until a prescribed stopping
criterion is satisfied. Rank-adaptive variants additionally enrich the
interpolation spaces when the current ranks are insufficient
\cite{oseledets2010ttcross,chertkov2023anovaTT}.\\
\\
If \(n_k\leq n\) and the intermediate TT ranks are bounded by \(r\), one sweep
typically requires
\begin{equation}
\mathcal{O}\!\left(dnr^2\right)
\end{equation}
oracle evaluations and
\begin{equation}
\mathcal{O}\!\left(dnr^3\right)
\end{equation}
arithmetic operations. In QTT form, where
\begin{equation}
L=\sum_{k=1}^{d}L_k
\end{equation}
is the number of binary modes, these estimates become
\(\mathcal{O}(Lr^2)\) and \(\mathcal{O}(Lr^3)\), respectively. These
complexities are conditional on the rank bound: they do not guarantee
dimension-independent performance because the required ranks may increase
with the complexity of the pricing tensor.\\
\\
Oracle calls are cached so that repeated requests for the same tensor entry do
not trigger additional pricing evaluations. Computational cost is therefore
measured by the number of distinct queried indices
\begin{equation}
N_{\mathrm{eval}}
=
\left|
\bigcup_s \Omega_s
\right|
\end{equation}
where \(\Omega_s\) denotes the set of indices requested during block update
\(s\).
Internal convergence is monitored through the relative change between
successive reconstructions
\begin{equation}
\varepsilon_{\mathrm{sweep}}^{(s)}
=
\frac{
\left\|
\widehat{\mathcal{A}}^{(s)}
-
\widehat{\mathcal{A}}^{(s-1)}
\right\|_F
}{
\left\|
\widehat{\mathcal{A}}^{(s)}
\right\|_F
}
\end{equation}
with $||.||_F$ the Frobenius norm. 
Accuracy is assessed independently on a set
\(\Omega_{\mathrm{val}}\) excluded from reconstruction
\begin{equation}
\varepsilon_{\mathrm{val}}
=
\left[
\frac{
\displaystyle
\sum_{\boldsymbol{i}\in\Omega_{\mathrm{val}}}
\left|
\mathcal{A}(\boldsymbol{i})
-
\widehat{\mathcal{A}}(\boldsymbol{i})
\right|^2
}{
\displaystyle
\sum_{\boldsymbol{i}\in\Omega_{\mathrm{val}}}
\left|
\mathcal{A}(\boldsymbol{i})
\right|^2
}
\right]^{1/2}
\end{equation}
The sweep criterion measures algorithmic stabilization and does not constitute
an accuracy bound on unseen tensor entries.\\
\\
When the pricing oracle is simulation-based, common random numbers and a fixed
random seed are used during reconstruction. TT-cross then operates on a fixed
random realization of the pricing tensor rather than on inconsistent values
generated at successive calls. Accuracy with respect to the underlying option
price is subsequently assessed using independent simulations.
\subsection{Off-Grid Interpolation and Gaussian-Process Interpretation}
\label{subsec:off-grid-interpolation}
Let $\boldsymbol{x}$ be a parameter vector that does not coincide with a grid node :
\begin{equation}
    \boldsymbol{x} = (x_1,\ldots,x_d) \in \Theta = \Theta_1\times\cdots\times\Theta_d
\end{equation}
The reconstructed
tensor $\widehat{\mathcal{A}}$ must then be combined with an interpolation
operator to define an off-grid pricing surrogate.
That's why, we mathematically define multilinear interpolation. 
For each coordinate $k$, let
\[
x_{k,j_k}
\leq
x_k
\leq
x_{k,j_k+1}
\]
denote the neighboring grid nodes and define
\begin{equation}
t_k(\boldsymbol{x}) = \frac{x_k-x_{k,j_k}} {x_{k,j_k+1}-x_{k,j_k}} \in[0,1]
\end{equation}
with interpolation weights $w_{k,0}(\boldsymbol{x}) = 1-t_k(\boldsymbol{x})$ and  $w_{k,1}(\boldsymbol{x}) = t_k(\boldsymbol{x})$.  Let $\boldsymbol{j} = (j_1,\ldots,j_d)$ denote the lower corner of the corresponding grid cell. The multilinear interpolant is
\begin{equation}
\widehat{V}_{\mathrm{TT}}(\boldsymbol{x}) = \sum_{\boldsymbol{b}\in\{0,1\}^d} \left[ \prod_{k=1}^{d} w_{k,b_k}(\boldsymbol{x}) \right] \widehat{\mathcal{A}} (j_1+b_1,\ldots,j_d+b_d)
\end{equation}
The weights are nonnegative and satisfy
\begin{equation}
\sum_{\boldsymbol{b}\in\{0,1\}^d} \prod_{k=1}^{d} w_{k,b_k}(\boldsymbol{x}) = 1
\end{equation}
so that the interpolant is a convex combination of the $2^d$ corner values. \\
\\
When $\widehat{\mathcal{A}}$ is directly available in TT form, interpolation can
also be incorporated into the tensor contraction. If $G^{(k)}(i_k) \in \mathbb{R}^{r_{k-1}\times r_k}$
denotes the $k$th TT core slice, define
\begin{equation}
B^{(k)}(x_k) = w_{k,0}(\boldsymbol{x})G^{(k)}(j_k) + w_{k,1}(\boldsymbol{x})G^{(k)}(j_k+1)
\end{equation}
Then
\begin{equation}
\widehat{V}_{\mathrm{TT}}(\boldsymbol{x}) = B^{(1)}(x_1) B^{(2)}(x_2) \cdots B^{(d)}(x_d)
\end{equation}
which avoids explicitly constructing all $2^d$ corner values. Equivalently, defining the rank-one interpolation tensor
\begin{equation}
\mathcal{W}(\boldsymbol{x}) = \boldsymbol{w}_1(x_1) \otimes\cdots\otimes \boldsymbol{w}_d(x_d)
\end{equation}
gives
\begin{equation}
\widehat{V}_{\mathrm{TT}}(\boldsymbol{x}) = \left\langle \widehat{\mathcal{A}}, \mathcal{W}(\boldsymbol{x}) \right\rangle
\end{equation}
Where $\langle .,.\rangle$ is the Frobenius inner product between two tensors of identical dimensions. In our implementation, the same interpolated value is obtained by evaluating
the reconstructed QTT at the relevant cell corners and applying the multilinear
weights.
Now, denoting $\mathcal{I}_h$ the multilinear interpolation operator associated
with the physical grid. The complete TT surrogate is
\begin{equation}
\widehat{V}_{\mathrm{TT}}(\boldsymbol{x}) = \mathcal{I}_h\widehat{\mathcal{A}}(\boldsymbol{x})
\end{equation}
Its error admits the decomposition
\begin{equation}
V(\boldsymbol{x})
-
\widehat{V}_{\mathrm{TT}}(\boldsymbol{x})
=
\underbrace{
V(\boldsymbol{x})
-
\mathcal{I}_h\mathcal{A}(\boldsymbol{x})
}_{\text{grid interpolation error}}
+
\underbrace{
\mathcal{I}_h
\left(
\mathcal{A}-\widehat{\mathcal{A}}
\right)(\boldsymbol{x})
}_{\text{propagated TT reconstruction error}}
\label{eq:deterministic-error-decomposition}
\end{equation}
This decomposition will be used throughout the remainder of the paper. It
separates the error imposed by the physical grid from the error introduced by
the tensor reconstruction.\\
\\
Since multilinear interpolation is a convex combination of nodal values, we can upper bound it this way 
\begin{equation}
\left|
\mathcal{I}_h
\left(
\mathcal{A}-\widehat{\mathcal{A}}
\right)(\boldsymbol{x})
\right|
\leq
\max_{\boldsymbol{b}\in\{0,1\}^d}
\left|
\mathcal{A}(\boldsymbol{j}+\boldsymbol{b})
-
\widehat{\mathcal{A}}(\boldsymbol{j}+\boldsymbol{b})
\right|
\end{equation}
Hence interpolation does not amplify the maximum nodal reconstruction error
inside a cell.\\
\\
Let
\begin{equation}
C_{\boldsymbol{j}} = \prod_{k=1}^{d} [x_{k,j_k},x_{k,j_k+1}]
\end{equation}
denote the cell containing $\boldsymbol{x}$ and define
\begin{equation}
h_{k,\boldsymbol{j}} = x_{k,j_k+1}-x_{k,j_k}
\end{equation}
If $V$ is twice continuously differentiable inside $C_{\boldsymbol{j}}$, a
standard local estimate gives
\[
\left|
V(\boldsymbol{x})
-
\mathcal{I}_h\mathcal{A}(\boldsymbol{x})
\right|
\leq
\frac{1}{8}
\sum_{k=1}^{d}
h_{k,\boldsymbol{j}}^2
\sup_{\boldsymbol{z}\in C_{\boldsymbol{j}}}
\left|
\frac{\partial^2 V}
{\partial x_k^2}(\boldsymbol{z})
\right|
\]
The interpolation error is therefore controlled jointly by the local grid
spacing and the directional curvature of the pricing surface. This observation
will motivate the adaptive grids introduced later.\\
\\
Multilinear interpolation can also be interpreted through a limiting
Gaussian-process construction. Consider first a zero-mean Gaussian process in
one dimension with Laplacian covariance kernel
\begin{equation}
k_{\ell}(x,x')
=
\exp\!\left(
-\frac{|x-x'|_1}{\ell}
\right),
\qquad
\ell>0
\end{equation}
where $\ell$ is the sole hyperparameter of the kernel and $|.|_1$ is the 1-Manhattan distance. 
Given noiseless observations
\(\boldsymbol{y}=(y_1,\ldots,y_N)^{\mathsf T}\) at distinct nodes
\(X=\{x_1,\ldots,x_N\}\), the posterior mean is
\begin{equation}
\mu_{\ell}(x)
=
\boldsymbol{k}_{\ell}(x,X)^{\mathsf T}
K_{\ell}(X,X)^{-1}
\boldsymbol{y}
\end{equation}
where
\begin{equation}
\boldsymbol{k}_{\ell}(x,X)
=
\begin{pmatrix}
k_{\ell}(x,x_1) &
\cdots &
k_{\ell}(x,x_N)
\end{pmatrix}^{\mathsf T}
\end{equation}
and the Gram matrix
\(K_{\ell}(X,X)\in\mathbb{R}^{N\times N}\) is defined entrywise by
\begin{equation}
\bigl[K_{\ell}(X,X)\bigr]_{ij}
=
k_{\ell}(x_i,x_j),
\qquad
i,j=1,\ldots,N
\end{equation}
For $x\in[x_j,x_{j+1}]$, the Markov structure of the exponential covariance
implies that the posterior mean depends only on the two neighboring
observations. Using
\begin{equation}
\sinh\!\left(\frac{z}{L}\right) = \frac{z}{L} + \mathcal{O}\!\left(L^{-3}\right), \qquad L\rightarrow\infty
\end{equation}
gives
\[
\mu_L(x)
\longrightarrow
y_j
\frac{x_{j+1}-x}
{x_{j+1}-x_j}
+
y_{j+1}
\frac{x-x_j}
{x_{j+1}-x_j}
\]
Thus, the large-length-scale limit of the noise-free Laplacian-GP posterior mean
is ordinary linear interpolation.\\
\\
In $d$ dimensions, the separable kernel is 
\begin{equation}
k_{\boldsymbol{L}}
(\boldsymbol{x},\boldsymbol{x}')
=
\prod_{k=1}^{d}
\exp\!\left(
-\frac{|x_k-x_k'|}{L_k}
\right)
\end{equation}
On a Cartesian product grid, the covariance structure factorizes across
coordinates, and the corresponding large-length-scale limit yields multilinear
interpolation.\\
\\
The strict limit $L_k\rightarrow\infty$ produces a degenerate constant
covariance kernel. The relevant statement is therefore a limit of posterior
means for finite length scales rather than a Gaussian process defined directly
at $L_k=\infty$. In the original STN-GPR experiments, marginal-likelihood
optimization drove the length scales toward large values, motivating
multilinear interpolation as the operational STN predictor.\\
\\
An exact Gaussian Process Regression model is also considered as an independent
benchmark for off-grid prediction. It provides a useful reference against which
the interpolation accuracy of the TT approximation can be assessed, but it is
not incorporated into the proposed reconstruction procedure.\\
\\
Multilinear interpolation is continuous but not generally continuously
differentiable across cell boundaries. In particular, its pure second
derivatives vanish within each cell and are not classically defined at cell
interfaces. It may therefore provide accurate price estimates while remaining
unsuitable for the direct computation of Gamma and other higher-order local
sensitivities.
\subsection{Conditions for Efficient Reconstruction}
\label{subsec:reconstruction-conditions}
The preceding construction separates nodal reconstruction from off-grid
interpolation. Its efficiency first requires the nodal pricing tensor to admit
a sufficiently accurate QTT representation with moderate ranks. For a
prescribed tolerance \(\varepsilon_{\mathrm{TT}}>0\), we assume that there
exists a tensor
\(\mathcal{A}_{\boldsymbol{r}}\in\operatorname{TT}(\boldsymbol{r})\) such that
\begin{equation}
\frac{
\left\|
\mathcal{A}-\mathcal{A}_{\boldsymbol{r}}
\right\|_F
}{
\left\|\mathcal{A}\right\|_F
}
\leq
\varepsilon_{\mathrm{TT}}
\end{equation}
where \(\operatorname{TT}(\boldsymbol{r})\) denotes the set of tensors whose TT
ranks do not exceed
\(\boldsymbol{r}=(r_0,\ldots,r_L)\). The required ranks depend on the coordinate
system, the physical grid mapping, and the ordering of the quantized modes.
Smoothness of the underlying pricing function alone does not guarantee a
low-rank QTT representation.\\
\\
Efficient reconstruction also requires stable interpolation pivots. Let \(P_k\)
denote the square pivot matrix selected from the \(k\)-th unfolding during the
cross procedure. Its conditioning is measured by
\begin{equation}
\kappa_2(P_k)
=
\left\|P_k\right\|_2
\left\|P_k^{-1}\right\|_2,
\qquad
k=1,\ldots,L-1
\end{equation}
Near-singular pivot matrices can amplify pricing, rounding, and truncation
errors even when the tensor is approximately low-rank. The \emph{maxvol}
selection reduces this risk but does not provide an unconditional global error
bound. \\
\\
Finally, the pricing oracle must assign a consistent value to each grid index.
This condition is immediate for deterministic pricing methods. For a
simulation-based oracle, the random numbers are held fixed throughout the
reconstruction, so that repeated evaluations define the same realization
\begin{equation}
\widetilde{\mathcal{A}}_{\mathrm{CRN}}(\boldsymbol{i})
=
\mathcal{A}(\boldsymbol{i})
+
\eta_{\mathrm{CRN}}(\boldsymbol{i})
\end{equation}
Independent simulations are subsequently used to assess accuracy with respect
to the underlying option value rather than to this fixed noisy realization.\\
\\
Under these conditions, increasing the TT-cross evaluation budget can improve
the approximation of the nodal tensor. It cannot eliminate errors inherited
from the pricing oracle or errors introduced when interpolating between grid
nodes. The latter depend on how well the chosen coordinates and physical grid
resolve the geometry of the pricing surface, which motivates the transformations
introduced in the following section.
\subsection{Limitations of Absolute-Strike Uniform Grids}
\label{subsec:absolute-strike-limitations}
Many option-pricing problems exhibit positive homogeneity with respect to the
underlying asset levels and the strike. Let
\begin{equation}
\boldsymbol{S} = (S_1,\ldots,S_{d_S})
\end{equation}
denote the underlying asset vector, where $d_S$ is the number of assets, and
let $\boldsymbol{\xi}$ collect the
remaining model parameters. For a positively homogeneous payoff of degree one,
the corresponding option value satisfies
\begin{equation}
V( \lambda S_1,\ldots,\lambda S_{d_S}, \lambda K, r,T,\boldsymbol{\xi} ) = \lambda V( S_1,\ldots,S_{d_S}, K,r,T,\boldsymbol{\xi} ), \qquad \lambda>0
\label{eq:homogeneity-preview}
\end{equation}
This property holds for both European expectations and American Snell envelopes,
since conditional expectation and maximization preserve positive homogeneity. \\
\\
Let $B(\boldsymbol{S})$ denote a positively homogeneous reference scale for the
underlying state, such as the spot level in the single-asset case or an
appropriate basket value in the multi-asset case. A natural dimensionless strike
coordinate is the log-moneyness
\begin{equation}
m = \log\!\left( \frac{K}{B(\boldsymbol{S})} \right)
\label{eq:basket-log-moneyness}
\end{equation}
The at-the-money region $K = B(\boldsymbol{S})$ is therefore mapped to the fixed coordinate \begin{equation}
m=0
\label{eq:atm-hyperplane}
\end{equation}
This distinction has direct consequences for grid resolution. On a uniform
strike grid with spacing $\Delta K$, the corresponding local log-moneyness
spacing near the money satisfies
\begin{equation}
\Delta m = \log\!\left( 1+\frac{\Delta K}{B(\boldsymbol{S})} \right) \approx \frac{\Delta K}{B(\boldsymbol{S})}
\end{equation}
Hence, the effective relative resolution varies with the basket level. A single
absolute-strike spacing cannot provide comparable moneyness resolution across
all spot configurations.\\
\\
Positive homogeneity also makes the scale dependence of strike curvature
explicit. Writing
\begin{equation}
V(\boldsymbol{S},K,\ldots) = B(\boldsymbol{S})\,v(m,\ldots)
\label{eq:normalized-price-preview}
\end{equation}
gives the chain-rule identities
\begin{equation}
\frac{\partial V}{\partial K}
=
e^{-m}\frac{\partial v}{\partial m}
\qquad
\frac{\partial^2V}{\partial K^2}
=
\frac{e^{-2m}}{B(\boldsymbol{S})}
\left(
\frac{\partial^2v}{\partial m^2}
-
\frac{\partial v}{\partial m}
\right)
\label{eq:normalized-strike-derivatives}
\end{equation}
Combined with the interpolation estimate established in
Section~\ref{subsec:off-grid-interpolation},
this relation shows that uniform absolute-strike spacing leads to
scale-dependent interpolation accuracy. The effect is particularly relevant
near the at-the-money region, where option convexity is typically concentrated. \\
\\
This also clarifies the role of the TT-cross sampling budget. The error
decomposition introduced above separates the grid interpolation error from the
nodal reconstruction error. Increasing the TT-cross budget acts only on the
latter. Once the reconstruction error becomes smaller than the interpolation
floor imposed by the physical grid, additional tensor evaluations or larger
ranks cannot compensate for an insufficiently resolved strike discretization. \\
\\
Absolute-strike coordinates may also be unfavorable for tensor compression. The
at-the-money ridge $K$
moves jointly with the spot coordinates and is generally oblique to the tensor
axes. Tensor unfoldings may consequently exhibit stronger apparent coupling
between the strike dimension and the underlying assets, potentially increasing
the required TT or QTT ranks. \\
\\
The log-moneyness transformation instead maps this ridge to
\eqref{eq:atm-hyperplane}
and separates the overall price scale from the relative strike position. This
does not guarantee low ranks, but it aligns the tensor coordinates more closely
with the financial geometry of the pricing function.\\
\\
For American options, the same issue is compounded by the early-exercise
frontier. The continuation and exercise regions meet along a state-dependent
surface and introduce localized changes in regularity. A coarse
absolute-strike grid may undersample this transition, while LSMC noise may
obscure it further. TT-cross may then either smooth over the transition or
increase its ranks in an attempt to reproduce noisy local variations.\\
\\
These limitations are not intrinsic to QTT. As established above, binary
tensorization constrains the index structure but does not require uniform
physical spacing. The same QTT architecture can therefore be combined with
nonlinear coordinate mappings and non-uniform grids. The limitations identified
here arise instead from the use of absolute-strike coordinates together with
uniform physical nodes. This motivates the homogeneity-based transformation and
convexity-guided grid constructions developed in the following sections.
\section{Homogeneity-Based Transformation and Log-Moneyness Coordinates}
\label{sec:homogeneity-transformation}
Section~\ref{subsec:absolute-strike-limitations} introduced basket
log-moneyness and identified the scale dependence created by an absolute-strike
grid. The present section provides the mathematical justification for that
change of coordinates. The argument has three steps. First, basket-option
prices inherit the homogeneity of the underlying dynamics and payoff. Second,
this homogeneity separates the common monetary scale from basket composition
and moneyness. Third, removing the resulting state-dependent translations can
improve the ranks of the discretized pricing tensor. The TT, QTT and
interpolation constructions themselves remain those of
Sections~\ref{subsec:tt-qtt}--\ref{subsec:off-grid-interpolation}.
\subsection{Homogeneity of Basket-Option Prices}
\label{subsec:price-homogeneity}
Let \(d_S\) denote the number of underlying assets and assume that their
risk-neutral dynamics satisfy the pathwise scaling property
\begin{equation}
\boldsymbol{S}_u^{t,\lambda\boldsymbol{s}}
=
\lambda\boldsymbol{S}_u^{t,\boldsymbol{s}}
\qquad
\lambda>0
\qquad
\mathbb{Q}\text{-a.s.}
\label{eq:pathwise-scaling}
\end{equation}
Correlated geometric Brownian motions satisfy
\eqref{eq:pathwise-scaling} \cite{black1973pricing,merton1973rational}, but the
argument below does not otherwise depend on lognormality. Let
\(B:(0,\infty)^{d_S}\to(0,\infty)\) be positively homogeneous of degree one and
define
\begin{equation}
\Phi_\eta(\boldsymbol{s},K)
=
\left[
\eta\left(B(\boldsymbol{s})-K\right)
\right]^+,
\qquad
\eta\in\{-1,1\}
\label{eq:basket-payoff}
\end{equation}
Weighted arithmetic baskets and weighted geometric baskets with weights
summing to one both satisfy the required homogeneity.

\begin{proposition}[Price homogeneity]
For both European and American exercise,
\begin{equation}
V
\left(
t,\lambda\boldsymbol{s},\lambda K,\boldsymbol{\vartheta}
\right)
=
\lambda
V
\left(
t,\boldsymbol{s},K,\boldsymbol{\vartheta}
\right)
\label{eq:price-homogeneity}
\end{equation}
for every \(\lambda>0\), provided that
\(\boldsymbol{\vartheta}\) contains only scale-free parameters.
\end{proposition}

\begin{proof}
Let \(\mathfrak{T}_{t,T}=\{T\}\) for a European option and let
\(\mathfrak{T}_{t,T}\) be the admissible stopping times for an American option
\cite{jaillet1990variational}.
Both prices can be written as
\begin{equation}
V(t,\boldsymbol{s},K,\boldsymbol{\vartheta})
=
\sup_{\tau\in\mathfrak{T}_{t,T}}
\mathbb{E}^{\mathbb{Q}}
\left[
e^{-r(\tau-t)}
\Phi_\eta
\left(
\boldsymbol{S}_\tau^{t,\boldsymbol{s}},K
\right)
\right]
\end{equation}
By \eqref{eq:pathwise-scaling} and the degree-one homogeneity of
\(\Phi_\eta\), the discounted payoff is multiplied by \(\lambda\) for every
admissible \(\tau\). Since \(\lambda>0\), expectation and supremum preserve this
factor, which proves \eqref{eq:price-homogeneity}.
\end{proof}

The same property is retained by the LSMC estimator used in the experiments
\cite{longstaff2001leastSquares,glasserman2004montecarlo}
when its regression state and basis functions are dimensionless and common
random numbers are used across scaled states. Conditional on the simulated
paths, exercise decisions are unchanged whereas cash flows are multiplied by
\(\lambda\). Remaining deviations from homogeneity are then numerical rather
than structural.

\subsection{Reduction to Relative Coordinates}
\label{subsec:relative-coordinates}
Equation~\eqref{eq:price-homogeneity} shows that the common monetary scale is a
redundant degree of freedom. Using the basket value as numeraire, define the
normalized composition
\begin{equation}
\overline{\boldsymbol{s}}
=
\frac{\boldsymbol{s}}{B(\boldsymbol{s})}
\label{eq:normalized-composition}
\end{equation}
which belongs to the cross-section
\begin{equation}
\mathcal{M}_B
=
\left\{
\boldsymbol{z}\in(0,\infty)^{d_S}
:
B(\boldsymbol{z})=1
\right\}
\label{eq:basket-cross-section}
\end{equation}
Under the regularity condition
\(\nabla B(\boldsymbol{z})\neq\boldsymbol{0}\),
this cross-section has dimension \(d_S-1\).\\
\\
We retain the log-moneyness \(m\) introduced in
\eqref{eq:basket-log-moneyness} and define a single normalized pricing
function
\begin{equation}
v
\left(
t,\overline{\boldsymbol{s}},m,\boldsymbol{\vartheta}
\right)
=
V
\left(
t,\overline{\boldsymbol{s}},e^m,\boldsymbol{\vartheta}
\right)
\label{eq:normalized-pricing-function}
\end{equation}

\begin{proposition}[Exact scale separation]
The original price admits the factorization
\begin{equation}
V
\left(
t,\boldsymbol{s},K,\boldsymbol{\vartheta}
\right)
=
B(\boldsymbol{s})
v
\left(
t,\overline{\boldsymbol{s}},m,\boldsymbol{\vartheta}
\right)
\label{eq:exact-scale-separation}
\end{equation}
\end{proposition}

This follows immediately from \eqref{eq:price-homogeneity} with
\(\lambda=B(\boldsymbol{s})^{-1}\) and
\(K/B(\boldsymbol{s})=e^m\).\\
\\
The factorization reduces the intrinsic dimension of
\((\boldsymbol{s},K)\) by one: the normalized composition contributes
\(d_S-1\) degrees of freedom and \(m\) contributes one. Basket composition
cannot generally be removed, because two normalized spot vectors with the same
basket value may induce different future basket distributions. It disappears
only in special cases, such as a geometric basket under lognormal dynamics
with fixed scale-free parameters.\\
\\
Relative to the original STN-GPR parameterization
\cite{Gribben2026STNGPR}, the numerical construction used in this work adopts
a less restrictive transformation than the complete quotient reduction. It
retains the absolute spot coordinates and replaces only the strike by
log-moneyness. Accordingly, the transformed surface reconstructed by the
surrogate is
\begin{equation}
\widetilde{V}
\left(
t,\boldsymbol{s},m,\boldsymbol{\vartheta}
\right)
=
V
\left(
t,\boldsymbol{s},B(\boldsymbol{s})e^m,\boldsymbol{\vartheta}
\right)
=
B(\boldsymbol{s})
v
\left(
t,\overline{\boldsymbol{s}},m,\boldsymbol{\vartheta}
\right)
\label{eq:transformed-pricing-surface}
\end{equation}
Thus, the exact quotient representation explains the financial structure,
whereas \(\widetilde V\) is the object actually sampled by TT-cross
\cite{oseledets2010ttcross}. This practical choice does not reduce the nominal
tensor order; it removes the redundant spot-dependent displacement of the
strike profile.
\subsection{Consequences of the Log-Moneyness Transformation}
\label{subsec:log-moneyness-consequences}
The coordinate and its at-the-money interpretation have already been
established in Section~\ref{subsec:absolute-strike-limitations}. Only the
consequences required by the numerical method are recorded here.\\
\\
For a log-moneyness grid
\(m_0<\cdots<m_{n_m-1}\), its strike nodes at a fixed spot state are
\begin{equation}
K_j(\boldsymbol{s})
=
B(\boldsymbol{s})e^{m_j}
\end{equation}
Hence,
\begin{equation}
\frac{
K_{j+1}(\boldsymbol{s})-K_j(\boldsymbol{s})
}{
K_j(\boldsymbol{s})
}
=
e^{m_{j+1}-m_j}-1
\label{eq:relative-strike-spacing}
\end{equation}
which is independent of the basket level. Equation
\eqref{eq:relative-strike-spacing} follows immediately by cancellation of
\(B(\boldsymbol{s})\); a separate proof is unnecessary.\\
\\
The strike-derivative identities are those already derived in
\eqref{eq:normalized-strike-derivatives}. One additional chain-rule correction
is needed when market Delta is computed while holding the absolute strike
fixed:
\begin{equation}
\left.
\frac{\partial V}{\partial s_i}
\right|_K
=
\left.
\frac{\partial\widetilde{V}}{\partial s_i}
\right|_m
-
\frac{
\partial_{s_i}B(\boldsymbol{s})
}{
B(\boldsymbol{s})
}
\frac{\partial\widetilde{V}}{\partial m}
\label{eq:market-delta-correction}
\end{equation}
This distinction becomes important when Greeks are extracted from a surrogate
trained in transformed coordinates.\\
\\
Likewise, the interpolation estimate of
Section~\ref{subsec:off-grid-interpolation} applies directly with the local
spacing \(h_j=m_{j+1}-m_j\). If the transformed price is twice continuously differentiable with respect to $m$ on this interval, linear interpolation satisfies
\[
\left|
\widetilde{V}
-
\mathcal{I}_m\widetilde{V}
\right|
\leq
\frac{h_j^2}{8}
\sup_{\xi\in[m_j,m_{j+1}]}
\left|
\frac{\partial^2\widetilde{V}}{\partial m^2}(\xi)
\right|
\]
Using $\widetilde{V}=B(\boldsymbol{s})u$ gives the scale-free estimate
\[
\frac{
\left|
\widetilde{V}
-
\mathcal{I}_m\widetilde{V}
\right|
}{
B(\boldsymbol{s})
}
\leq
\frac{h_j^2}{8}
\sup_{\xi\in[m_j,m_{j+1}]}
\left|
\frac{\partial^2u}{\partial m^2}(\xi)
\right|
\]
The normalized interpolation error is thus controlled by the log-moneyness spacing and the curvature of a dimensionless pricing function, rather than by an absolute strike increment.\\
\\
If the original market domain uses
\(K\in[K_{\min},K_{\max}]\) and
\(\boldsymbol{s}\in\mathcal{S}\), complete coverage is obtained from the
rectangular envelope
\begin{equation}
m_{\min}
=
\log\left(
\frac{K_{\min}}{B_{\max}}
\right),
\qquad
m_{\max}
=
\log\left(
\frac{K_{\max}}{B_{\min}}
\right)
\label{eq:moneyness-domain}
\end{equation}
where
\[
B_{\min}
=
\min_{\boldsymbol{s}\in\mathcal{S}}B(\boldsymbol{s}),
\qquad
B_{\max}
=
\max_{\boldsymbol{s}\in\mathcal{S}}B(\boldsymbol{s})
\]
This envelope may contain inverse images outside the original strike interval
for some spot states, but it contains the image of every state in the original
domain.\\
\\
Finally, QTT compatibility requires only the power-of-two cardinality of the
index set, as established in Section~\ref{subsec:tt-qtt}; it does not require
uniform physical spacing \cite{khoromskij2011qtt,oseledets2011waveletQTT}. No
binary modes or tensorization rules are therefore redefined here. The nonlinear
allocation of the \(m\)-nodes is developed in Section~4.
\subsection{Impact on Tensor Compressibility}
\label{subsec:compressibility-impact}
The exact TT ranks and their unfolding interpretation were defined in
\eqref{eq:tt-unfolding-rank} \cite{oseledets2011tensorTrain}. For a tolerance
\(\varepsilon>0\), we use the corresponding numerical rank
\begin{equation}
r_k^{(\varepsilon)}
=
\min
\left\{
r:
\sum_{j>r}
\left(
\sigma_j^{(k)}
\right)^2
\leq
\varepsilon^2
\left\|
\mathcal{A}^{\langle k\rangle}
\right\|_F^2
\right\}
\label{eq:numerical-tt-rank}
\end{equation}
where \(\sigma_j^{(k)}\) are the singular values of the \(k\)-th unfolding.
This extends the exact notion from Section~\ref{subsec:tt-qtt} without
redefining the tensor or its matricizations.\\
\\
To isolate the effect of the coordinate choice, first suppose that basket
composition and all other scale-free parameters are fixed, so that the
normalized price is a univariate profile \(f(m)\). For spot states
\(\boldsymbol{s}^{(i)}\), set
\(b_i=B(\boldsymbol{s}^{(i)})\). On a common moneyness grid,
\begin{equation}
A_{ij}^{\mathrm{rel}}
=
b_i f(m_j)
\label{eq:relative-rank-one-matrix}
\end{equation}
This matrix is the outer product of
\((b_i)_i\) and \((f(m_j))_j\), so its rank is one whenever the price is not
identically zero. This observation is immediate from
\eqref{eq:relative-rank-one-matrix} and requires no separate proof. \\
\\
On an absolute log-strike grid, let
\(\kappa_j=\log K_j\), \(\beta_i=\log b_i\),
\(\delta_i=\beta_i-\beta_0\), and
\(\overline m_j=\kappa_j-\beta_0\). The corresponding matrix is
\begin{equation}
A_{ij}^{\mathrm{abs}}
=
b_i f(\overline m_j-\delta_i)
\label{eq:absolute-translated-matrix}
\end{equation}
Its rows are translated copies of the same profile. A Taylor expansion gives
\begin{equation}
A_{ij}^{\mathrm{abs}}
=
\sum_{p=0}^{P}
\left[
b_i\frac{(-\delta_i)^p}{p!}
\right]
f^{(p)}(\overline m_j)
+
b_iR_{P+1}(i,j)
\label{eq:translation-separation}
\end{equation}
The truncated sum has rank at most \(P+1\). If
\(\lvert\delta_i\rvert\leq D\) and
\(\sup_{\mathcal{I}_D}\lvert f^{(P+1)}\rvert\leq M_{P+1}\), Taylor's theorem
yields
\begin{equation}
\left\|
A^{\mathrm{abs}}
-
A_P^{\mathrm{abs}}
\right\|_F
\leq
b_{\max}\sqrt{n_Sn_K}
\frac{
M_{P+1}D^{P+1}
}{
(P+1)!
}
\label{eq:absolute-rank-bound}
\end{equation}
The basket-scale range \(D\) and the derivatives of the moneyness profile
therefore control the additional separation rank required in absolute
coordinates. The relative transformation removes this translation parameter
exactly.\\
\\
For the general normalized surface in
\eqref{eq:exact-scale-separation}, suppose that
\begin{equation}
v
\left(
\overline{\boldsymbol{s}},m,\boldsymbol{\vartheta}
\right)
\approx
\sum_{\ell=1}^{R}
a_\ell
\left(
\overline{\boldsymbol{s}}
\right)
b_\ell
\left(
m,\boldsymbol{\vartheta}
\right)
\label{eq:separated-normalized-surface}
\end{equation}
at tolerance \(\varepsilon\). After discretization, each term is an outer
product across the unfolding that separates the spot modes from the remaining
modes. Row scaling by \(B(\boldsymbol{s})\) does not increase its rank.
Consequently, the corresponding numerical rank is at most \(R\). This is a
direct consequence of \eqref{eq:separated-normalized-surface}, so an additional
formal proof would add no information.\\
\\
No universal inequality between QTT ranks in the absolute and relative
coordinates follows. The transformation removes one
identified source of rank growth, but the remaining ranks still depend on
basket composition, mode ordering, heterogeneous dynamics, early exercise and
pricing-oracle noise. Their effect on storage and TT-cross cost follows from
the complexity relations already established in
Sections~\ref{subsec:tt-qtt} and~\ref{subsec:tt-cross}. The claim is therefore
tested empirically at fixed tensor shape, evaluation budget and reconstruction
tolerance. Section~4 now complements this algebraic reorganization by placing
more physical nodes where the normalized moneyness surface is most difficult
to interpolate.

\section{Convexity-guided adaptive QTT grids}
\subsection{Convexity around the At-the-Money Region}
Section~\ref{subsec:absolute-strike-limitations} and section~\ref{subsec:relative-coordinates} established the financial role of log-moneyness and the associated separation between monetary scale and relative option geometry. We now identify the region in which a fixed number of QTT nodes should provide the highest physical resolution. Let the basket function be either the weighted arithmetic average
\begin{equation}
B_{\mathrm{A}}(\boldsymbol{s})
=
\sum_{i=1}^{d_S} w_i s_i
\end{equation}
or the weighted geometric average
\begin{equation}
B_{\mathrm{G}}(\boldsymbol{s})
=
\prod_{i=1}^{d_S} s_i^{w_i}
\end{equation}
where \(w_i>0\) and
\begin{equation}
\sum_{i=1}^{d_S}w_i=1
\end{equation}
Both basket functions are positively homogeneous:
\begin{equation}
B_{\mathrm{A}}(\lambda\boldsymbol{s})
=
\lambda B_{\mathrm{A}}(\boldsymbol{s})
\qquad
B_{\mathrm{G}}(\lambda\boldsymbol{s})
=
\lambda B_{\mathrm{G}}(\boldsymbol{s})
\end{equation}
for every \(\lambda>0\).
\begin{proposition}[Strike convexity for arithmetic and geometric basket puts]
Let \(B\) denote either \(B_{\mathrm{A}}\) or \(B_{\mathrm{G}}\). The European and American basket-put prices
\begin{equation}
V_B^{\mathrm{E}}(t,\boldsymbol{s},K)
=
\mathbb{E}_t
\left[
e^{-r(T-t)}
\left(
K-B(\boldsymbol{S}_T)
\right)^+
\right]
\end{equation}
and
\begin{equation}
V_B^{\mathrm{A}}(t,\boldsymbol{s},K)
=
\sup_{\tau\in\mathcal{T}_{t,T}}
\mathbb{E}_t
\left[
e^{-r(\tau-t)}
\left(
K-B(\boldsymbol{S}_\tau)
\right)^+
\right]
\end{equation}
are convex functions of \(K\).
\end{proposition}

\begin{proof}
For every fixed basket value \(x\geq 0\), the function
\begin{equation}
K
\longmapsto
(K-x)^+
=
\max\{0,K-x\}
\end{equation}
is the pointwise maximum of two affine functions and is therefore convex. Multiplication by a positive discount factor and conditional expectation preserve convexity. This proves the result for the European price and for the discounted expected payoff associated with every fixed stopping time \(\tau\). \\
\\
The American price is the pointwise supremum over all admissible stopping times of these convex functions. Since a pointwise supremum of convex functions is convex, the result also holds for the American price.
\end{proof}
Consequently, both arithmetic and geometric basket-put prices satisfy
\begin{equation}
\frac{\partial^2 V_B^{\mathrm{E}}}{\partial K^2}
\geq 0
\qquad
\frac{\partial^2 V_B^{\mathrm{A}}}{\partial K^2}
\geq 0
\end{equation}
in the classical sense wherever the derivatives exist and in the distributional sense otherwise.
For the European option, the convexity admits a direct probabilistic interpretation. Let \(F_{B,T}(\,\cdot\mid\boldsymbol{s})\) denote the conditional distribution function of \(B(\boldsymbol{S}_T)\). Then
\begin{equation}
\frac{\partial V_B^{\mathrm{E}}}{\partial K}
=
e^{-r(T-t)}
F_{B,T}(K\mid\boldsymbol{s})
\end{equation}
If the conditional law admits a density \(f_{B,T}(\,\cdot\mid\boldsymbol{s})\), it follows that
\begin{equation}
\frac{\partial^2 V_B^{\mathrm{E}}}{\partial K^2}
=
e^{-r(T-t)}
f_{B,T}(K\mid\boldsymbol{s})
\geq 0
\end{equation}
This identity applies in particular to the arithmetic basket, even though the distribution of a weighted sum of correlated lognormal variables does not generally admit a closed-form expression. \\
\\
The homogeneity result established in Section 3.1 further implies
\begin{equation}
\frac{\partial^2 V_B}{\partial K^2}
(t,\lambda\boldsymbol{s},\lambda K,\boldsymbol{\vartheta})
=
\frac{1}{\lambda}
\frac{\partial^2 V_B}{\partial K^2}
(t,\boldsymbol{s},K,\boldsymbol{\vartheta})
\end{equation}
The magnitude of the strike curvature therefore depends on the monetary scale. To remove this dependence, set
\begin{equation}
B_0=B(\boldsymbol{s})
\qquad
K=B_0e^m
\end{equation}

\begin{definition}[Dimensionless strike convexity]
The dimensionless strike-convexity indicator is defined by
\begin{equation}
\chi_B(\boldsymbol{s},m,T,\boldsymbol{\vartheta})
=
\frac{K^2}{B_0}
\frac{\partial^2V_B}{\partial K^2}
(\boldsymbol{s},K,T,\boldsymbol{\vartheta})
\end{equation}
\end{definition}
The scaling identity gives
\begin{equation}
\chi_B
(
\lambda\boldsymbol{s},m,T,\boldsymbol{\vartheta}
)
=
\chi_B
(
\boldsymbol{s},m,T,\boldsymbol{\vartheta}
)
\end{equation}
Moreover, if
\begin{equation}
V_B(\boldsymbol{s},K,T,\boldsymbol{\vartheta})
=
B_0
v_B(m,T,\boldsymbol{\vartheta})
\end{equation}
then
\begin{equation}
\chi_B
=
\frac{\partial^2v_B}{\partial m^2}
-
\frac{\partial v_B}{\partial m}
\end{equation}
The indicator therefore removes both the arbitrary monetary scale and the first-order contribution generated by the exponential transformation \(K=B_0e^m\).
\paragraph{Short-maturity behaviour of the arithmetic basket}
Consider the arithmetic basket
\begin{equation}
B_{\mathrm{A},0}
=
\sum_{i=1}^{d_S}w_is_i
\end{equation}
and define the state-dependent normalized basket weights
\begin{equation}
\alpha_i(\boldsymbol{s})
=
\frac{w_is_i}{B_{\mathrm{A},0}}
\end{equation}
Under correlated geometric Brownian dynamics, let
\begin{equation}
\Sigma_{ij}
=
\sigma_i\sigma_j\rho_{ij}
\end{equation}
The instantaneous effective variance of the normalized arithmetic basket is
\begin{equation}
\sigma_{\mathrm{A}}^2(\boldsymbol{s})
=
\boldsymbol{\alpha}(\boldsymbol{s})^{\mathsf T}
\Sigma
\boldsymbol{\alpha}(\boldsymbol{s})
\end{equation}
Its instantaneous drift is
\begin{equation}
\mu_{\mathrm{A}}(\boldsymbol{s})
=
\sum_{i=1}^{d_S}
\alpha_i(\boldsymbol{s})(r-q_i)
\end{equation}
For a short remaining maturity \(T\), a first-order stochastic expansion gives
\begin{equation}
\frac{
B_{\mathrm{A}}(\boldsymbol{S}_T)-B_{\mathrm{A},0}
}{
B_{\mathrm{A},0}
}
=
\mu_{\mathrm{A}}(\boldsymbol{s})T
+
\sigma_{\mathrm{A}}(\boldsymbol{s})\sqrt{T}\,Z
+
o_{\mathbb{P}}(\sqrt{T})
\end{equation}
where \(Z\sim\mathcal{N}(0,1)\). The transition of the put payoff is therefore smoothed over a log-moneyness region whose characteristic width is
\begin{equation}
O
\left(
\sigma_{\mathrm{A}}(\boldsymbol{s})\sqrt{T}
\right)
\end{equation}
The centre of this region satisfies
\begin{equation}
m_{\mathrm{A}}^\star(T)
=
O(T)
\end{equation}
while the maximum dimensionless convexity is of order
\begin{equation}
\chi_{\mathrm{A}}
\left(
m_{\mathrm{A}}^\star(T),T
\right)
=
O
\left(
\frac{1}{
\sigma_{\mathrm{A}}(\boldsymbol{s})\sqrt{T}
}
\right)
\end{equation}
Thus, the arithmetic-basket convexity concentrates around \(m=0\) and becomes increasingly localized as maturity approaches zero. This is the relevant behaviour for the arithmetic-basket American option considered in the numerical experiments.
\paragraph{Analytical geometric-basket benchmark}
Unlike the arithmetic basket, the geometric basket is lognormal under correlated geometric Brownian dynamics. It therefore provides an analytical benchmark through which the location and width of the convexity region can be computed exactly.\\
\\
Let
\begin{equation}
\sigma_{\mathrm{G}}^2
=
\boldsymbol{w}^{\mathsf T}
\Sigma
\boldsymbol{w}
\end{equation}
and
\begin{equation}
\mu_{\mathrm{G}}
=
r
-
\boldsymbol{w}^{\mathsf T}\boldsymbol{q}
-
\frac{1}{2}
\sum_{i=1}^{d_S}w_i\sigma_i^2
+
\frac{1}{2}\sigma_{\mathrm{G}}^2
\end{equation}
The terminal geometric basket satisfies
\begin{equation}
B_{\mathrm{G}}(\boldsymbol{S}_T)
=
B_{\mathrm{G},0}
\exp
\left[
\left(
\mu_{\mathrm{G}}
-
\frac{1}{2}\sigma_{\mathrm{G}}^2
\right)T
+
\sigma_{\mathrm{G}}\sqrt{T}\,Z
\right]
\end{equation}
where \(Z\sim\mathcal{N}(0,1)\). The normalized European put price is
\begin{equation}
v_{\mathrm{G}}^{\mathrm{P}}(m,T)
=
e^{m-rT}\Phi(-d_2)
-
e^{(\mu_{\mathrm{G}}-r)T}\Phi(-d_1)
\end{equation}

with

\begin{equation}
d_1
=
\frac{
-m+
\left(
\mu_{\mathrm{G}}
+
\frac{1}{2}\sigma_{\mathrm{G}}^2
\right)T
}{
\sigma_{\mathrm{G}}\sqrt{T}
}
\qquad
d_2
=
d_1-\sigma_{\mathrm{G}}\sqrt{T}
\end{equation}
Differentiation with respect to the strike gives
\begin{equation}
\frac{\partial^2V_{\mathrm{G}}^{\mathrm{E}}}{\partial K^2}
=
\frac{
e^{-rT}\varphi(d_2)
}{
K\sigma_{\mathrm{G}}\sqrt{T}
}
\end{equation}
and therefore
\begin{equation}
\chi_{\mathrm{G}}(m,T)
=
\frac{
e^{m-rT}\varphi(d_2)
}{
\sigma_{\mathrm{G}}\sqrt{T}
}
\end{equation}

\begin{lemma}[Location of the geometric-basket convexity ridge]
For every \(T>0\), the function
\begin{equation}
m
\longmapsto
\chi_{\mathrm{G}}(m,T)
\end{equation}
has a unique maximum at
\begin{equation}
m_{\mathrm{G}}^\star(T)
=
\left(
\mu_{\mathrm{G}}
+
\frac{1}{2}\sigma_{\mathrm{G}}^2
\right)T
\end{equation}
\end{lemma}

\begin{proof}
Direct differentiation gives
\begin{equation}
\frac{\partial}{\partial m}
\log\chi_{\mathrm{G}}
=
1+
\frac{d_2}{
\sigma_{\mathrm{G}}\sqrt{T}
}
\end{equation}
and
\begin{equation}
\frac{\partial^2}{\partial m^2}
\log\chi_{\mathrm{G}}
=
-
\frac{1}{
\sigma_{\mathrm{G}}^2T
}
<0
\end{equation}
The first-order condition gives
\begin{equation}
d_2
=
-\sigma_{\mathrm{G}}\sqrt{T}
\end{equation}
which is equivalent to the stated expression for \(m_{\mathrm{G}}^\star(T)\). Strict negativity of the second derivative proves uniqueness.
\end{proof}
As \(T\) approaches zero,
\begin{equation}
m_{\mathrm{G}}^\star(T)
=
O(T)
\longrightarrow 0
\end{equation}
and
\begin{equation}
\chi_{\mathrm{G}}
\left(
m_{\mathrm{G}}^\star(T),T
\right)
\sim
\frac{1}{
\sigma_{\mathrm{G}}\sqrt{2\pi T}
}
\end{equation}
The characteristic width of the ridge is
\begin{equation}
O
\left(
\sigma_{\mathrm{G}}\sqrt{T}
\right)
\end{equation}
The arithmetic and geometric baskets therefore exhibit the same leading short-maturity geometry: the convexity is concentrated around the at-the-money region, its width decreases as \(O(\sqrt{T})\), and its magnitude increases as \(O(T^{-1/2})\). The geometric benchmark is used in Figure 1 because its convexity field and ridge location are available analytically, whereas the arithmetic-basket argument does not require a closed-form distribution.\\
\\
Figure~\ref{fig:convexity-grid-comparison} illustrates the region in which the fixed QTT resolution should be concentrated. The white markers are included only to preview the physical nodes associated with the uniform and transformed discretizations. Their placement is not defined by the present convexity analysis. The construction and quantitative justification of the moneyness nodes are given in Section 4.2, while the placement of the maturity nodes is introduced separately in Section 4.3.\\
\\
For the five-asset equal-weight geometric basket used in Figure 1,
\begin{equation}
d_S=5
\qquad
w_i=\frac{1}{5}
\qquad
\sigma_i=20\%
\qquad
\rho_{ij}=0.3
\end{equation}
and
\begin{equation}
q_i=0
\qquad
r=3\%
\qquad
\sigma_{\mathrm{G}}=0.132665
\qquad
\mu_{\mathrm{G}}=0.0188
\end{equation}
Both panels use \(64\) moneyness nodes and \(8\) maturity nodes. At this stage, the analytical field identifies the financially relevant region but does not prescribe the coordinate transformation used to place those nodes.
\begin{figure}[htbp]
    \centering
    \includegraphics[
        width=\textwidth
    ]{figures/convexity-grid-comparison.png}
    \caption{
    Dimensionless strike convexity
    \(
    \chi_{\mathrm{G}}
    =
    (K^2/B_{\mathrm{G},0})V_{KK}
    \)
    for the European geometric-basket put. The cyan dashed curve is the analytical ridge
    \(
    m_{\mathrm{G}}^\star(T)
    \).
    The white circles preview the physical QTT nodes associated with the uniform and transformed grids. Their placement is not derived in the present subsection: the moneyness-node construction is introduced in Section 4.2 and the maturity-node construction in Section 4.3. The figure shows the short-maturity at-the-money region toward which the fixed physical resolution will be reallocated.
    }
    \label{fig:convexity-grid-comparison}
\end{figure}
The heat map provides a geometric interpretation of the adaptive construction. A uniform grid distributes its approximation capacity independently of the local structure of the pricing function. The convexity-guided grid instead preserves the same power-of-two cardinalities while moving physical nodes toward the region where the strike curvature is largest. The precise nonlinear mappings used to perform this redistribution are introduced in the following subsections.
\subsection{Hyperbolic-sine moneyness discretization}
Figure~\ref{fig:convexity-grid-comparison} identified the short-maturity at-the-money region as the region in which the pricing surface requires the highest physical resolution. We now define the placement of the moneyness nodes shown in that figure.\\
\\
The complete physical grid is a Cartesian product of the individual coordinate grids:
\begin{equation}
\mathcal{G}
=
\mathcal{G}_{S_1}
\times\cdots\times
\mathcal{G}_{S_{d_S}}
\times
\mathcal{G}_m
\times
\mathcal{G}_r
\times
\mathcal{G}_T
\end{equation}
and is therefore a separable metric space (see Appendix~\ref{app:separability-proof}). The componentwise grid maps preserve its Cartesian-product structure, so the discrete pricing array belongs naturally to the tensor-product space 
\[
\bigotimes_{k=1}^{d}\mathbb{R}^{n_k}
\]
For the five-asset experiments considered in this work, its shape is
\begin{equation}
32^5\times64\times8\times8
\end{equation}
The spot and interest-rate coordinates retain the uniform physical discretization of the baseline framework. The proposed adaptation modifies only the log-moneyness and maturity axes. The moneyness transformation is introduced in the present subsection, while the maturity transformation is defined in Section~\ref{subsec:short-maturity refinement}.\\
\\
The white markers in Figure~\ref{fig:convexity-grid-comparison} represent the two-dimensional projection
\begin{equation}
\mathcal{G}_m\times\mathcal{G}_T
\end{equation}
of this Cartesian grid. They show the physical nodes available to the tensor discretization and not the subset of multidimensional indices subsequently selected by TT-cross. The horizontal placement of the nodes is determined by the moneyness transformation defined below. Their vertical placement is explained separately in the next section.\\
\\
Now, let
\begin{equation}
m_{\min}<0<m_{\max}
\end{equation}
denote the prescribed bounds of the log-moneyness domain. We introduce the positive tail amplitudes
\begin{equation}
a_-=-m_{\min}
\qquad
a_+=m_{\max}
\end{equation}
and distinguish the physical coordinate \(m\) from a uniformly discretized computational coordinate
\begin{equation}
u\in[-1,1]
\end{equation}
The purpose of the coordinate map is to preserve the complete log-moneyness domain while allocating a larger fraction of the available physical resolution to a neighbourhood of \(m=0\). This is achieved without changing the number of QTT modes or their binary structure.
\begin{definition}[Hyperbolic-sine moneyness map]
For a concentration parameter \(\gamma>0\), define
\[
\phi_\gamma(u)
=
\begin{cases}
\displaystyle
a_-\frac{\sinh(\gamma u)}{\sinh(\gamma)},
& -1\leq u\leq 0
\\[1.2ex]
\displaystyle
a_+\frac{\sinh(\gamma u)}{\sinh(\gamma)},
& 0\leq u\leq 1
\end{cases}
\]
Given \(n_m=2^{L_m}\) computational nodes
\[
u_j=-1+\frac{2j}{n_m-1},
\qquad
j=0,\ldots,n_m-1
\]
the corresponding physical log-moneyness grid is
\[
m_j=\phi_\gamma(u_j)
\]
\end{definition}
The two branches are scaled independently because the admissible log-moneyness interval need not be symmetric around zero. Nevertheless, both branches share the same concentration parameter and therefore apply the same relative deformation to the in-the-money and out-of-the-money regions.
\begin{lemma}[Admissibility of the transformed grid]
The map \(\phi_\gamma\) is continuous and strictly increasing on \([-1,1]\), with
\[
\phi_\gamma(-1)=m_{\min},
\qquad
\phi_\gamma(0)=0,
\qquad
\phi_\gamma(1)=m_{\max}
\]
Consequently, the sequence \((m_j)_{j=0}^{n_m-1}\) is strictly increasing and covers the entire prescribed domain.
\end{lemma}

\begin{proof}
Continuity follows from the equality of the two branches at \(u=0\). Away from the origin,
\[
\phi_\gamma'(u)
=
\begin{cases}
\displaystyle
a_-\gamma\frac{\cosh(\gamma u)}{\sinh(\gamma)},
& u<0
\\[1.2ex]
\displaystyle
a_+\gamma\frac{\cosh(\gamma u)}{\sinh(\gamma)},
& u>0
\end{cases}
\]
which is strictly positive because \(a_->0\), \(a_+>0\), and \(\gamma>0\). The endpoint identities follow directly from the oddness of the hyperbolic sine.
\end{proof}
When \(a_-\neq a_+\), the one-sided derivatives of \(\phi_\gamma\) at the origin are different. The map is therefore continuous and piecewise smooth but not globally differentiable at \(u=0\). This does not introduce a singularity into the pricing model: \(\phi_\gamma\) is only an index-to-coordinate transformation used to construct the physical interpolation nodes.\\
\\
The local spacing is governed by \(\phi_\gamma'\). Since \(\cosh(\gamma u)\) is minimal at the origin and increases with \(\lvert u\rvert\), the smallest cells are located around the at-the-money region, whereas progressively wider cells are used in the tails. The inverse transformation on either side of the origin satisfies
\[
\lvert u\rvert
=
\frac{1}{\gamma}
\operatorname{arsinh}
\left(
\frac{\sinh(\gamma)}{a_\pm}\lvert m\rvert
\right)
\]
where \(a_-\) applies for \(m<0\) and \(a_+\) for \(m>0\).\\
\\
The associated continuous approximation of the physical node density is therefore
\[
\rho_\pm(m)
=
\frac{n_m-1}{2}
\left\lvert\frac{\mathrm{d}u}{\mathrm{d}m}\right\rvert
=
\frac{n_m-1}{2}
\frac{\sinh(\gamma)}{\gamma a_\pm}
\left[
1+
\left(
\frac{\sinh(\gamma)}{a_\pm}\lvert m\rvert
\right)^2
\right]^{-1/2}
\]
This density is maximal at \(m=0\) and decreases monotonically towards either endpoint. Hence, \(\gamma\) controls the strength of the spatial concentration without modifying the domain bounds or the number of nodes.\\
\\
Because \(n_m=2^{L_m}\) is even, zero is not itself a computational node. Instead, the two central nodes satisfy
\[
u_{n_m/2-1}=-\frac{1}{n_m-1},
\qquad
u_{n_m/2}=\frac{1}{n_m-1}
\]
and define the cell that straddles the ATM point.
\begin{proposition}[ATM refinement factor]
Let \(h_{\mathrm{ATM}}\) denote the width of the cell containing \(m=0\). Then
\[
h_{\mathrm{ATM}}
=
(a_-+a_+)
\frac{\sinh\left(\gamma/(n_m-1)\right)}
     {\sinh(\gamma)}
\]
For a uniform physical grid with the same bounds and cardinality,
\[
h_{\mathrm{unif}}
=
\frac{a_-+a_+}{n_m-1}
\]
and therefore
\[
\frac{h_{\mathrm{ATM}}}{h_{\mathrm{unif}}}
=
(n_m-1)
\frac{\sinh\left(\gamma/(n_m-1)\right)}
     {\sinh(\gamma)}
\]
Moreover, as \(n_m\rightarrow\infty\),
\[
\frac{h_{\mathrm{ATM}}}{h_{\mathrm{unif}}}
\longrightarrow
\frac{\gamma}{\sinh(\gamma)}
\]
\end{proposition}
For the discretization used in our experiments,
\[
n_m=64,
\qquad
\gamma=3
\]
which gives
\[
\frac{h_{\mathrm{ATM}}}{h_{\mathrm{unif}}}
\simeq 0.2996
\]
Thus, the ATM-straddling cell is approximately \(3.34\) times finer than its uniform-grid counterpart, while the number of tensor entries remains unchanged.\\
\\
This refinement has a direct interpolation interpretation. Let \(v(m)\) be twice continuously differentiable on a cell
\[
I_j=[m_j,m_{j+1}]
\]
and let \(\mathcal{I}_jv\) denote its linear interpolant. The classical local error bound is
\[
\sup_{m\in I_j}
\left|v(m)-\mathcal{I}_jv(m)\right|
\leq
\frac{h_j^2}{8}
\sup_{\xi\in I_j}
\left|v''(\xi)\right|
\]
where
\[
h_j=m_{j+1}-m_j
\]
The hyperbolic-sine map reduces \(h_j^2\) precisely in the region where the option-price curvature is largest. It therefore approximates a curvature-adapted discretization without requiring the grid to be recalibrated for every volatility, maturity, or correlation specification.\\
\\
The transformation is also compatible with QTT tensorization. The physical index \(j\) retains its binary representation
\[
j
=
\sum_{\ell=1}^{L_m}
2^{L_m-\ell}b_\ell,
\qquad
b_\ell\in\{0,1\}
\]
while only the deterministic map from \(j\) to \(m_j\) is modified. Hence, the QTT order, binary mode sizes, TT-cross sampling mechanism, and storage complexity remain unchanged \cite{khoromskij2011qtt,oseledets2011waveletQTT}. The grid deformation reallocates physical resolution without increasing the nominal tensor dimension.\\
\\
Figure~\ref{fig:sinh-moneyness-discretization} illustrates both effects. The first panel shows the nonlinear transformation from uniform computational indices to physical log-moneyness nodes. The second panel reports the resulting cell widths and makes the ATM refinement explicit.
\begin{figure}[t]
    \centering
    \includegraphics[width=\textwidth]{figures/sinh-moneyness-discretization.png}
    \caption{
    Hyperbolic-sine discretization of the log-moneyness coordinate.
    Panel (a) compares the uniform physical map with the proposed
    hyperbolic-sine map for \(n_m=64\) and \(\gamma=3\).
    Panel (b) shows the corresponding local cell widths.
    The vertical red line identifies the at-the-money point.
    The adaptive grid achieves approximately \(3.34\) times finer local
    resolution around the ATM region without increasing the QTT cardinality.
    }
    \label{fig:sinh-moneyness-discretization}
\end{figure}
The construction does not claim to solve an exact curvature-equidistribution problem. It provides instead a parsimonious and model-independent approximation to such a grid. Increasing \(\gamma\) improves ATM resolution but enlarges the tail cells, so the concentration parameter must ultimately be assessed through out-of-sample error, tail robustness, and exercise-frontier diagnostics. The ablation protocol developed in the following sections is designed to verify that the reduction of central interpolation error is not obtained at the cost of unacceptable deterioration in the tails.
\subsection{Short-maturity refinement}
\label{subsec:short-maturity refinement}
The log-moneyness transformation concentrates the payoff transition around the fixed hyperplane \(m=0\), but the spatial extent of this transition remains maturity-dependent. Under a diffusion model, the terminal payoff kink is smoothed over a region whose characteristic width is of order
\[
\sigma_{\mathrm{eff}}\sqrt{T}
\]
Consequently, the price surface develops a narrow convexity ridge as \(T\) approaches zero. At the money, the time value is typically of order \(\sqrt{T}\), while its maturity derivative may grow as \(T^{-1/2}\). Uniform maturity spacing therefore allocates insufficient resolution to the region where the surface varies most rapidly.\\
\\
Let
\[
0<T_{\min}<T_{\max}
\]
be the maturity bounds and let \(n_T=2^{L_T}\) be the number of maturity nodes. We introduce a uniform computational coordinate
\[
v_j=\frac{j}{n_T-1},
\qquad
j=0,\ldots,n_T-1
\]
and map it to physical maturity through
\[
T_j
=
T_{\min}
+
\left(T_{\max}-T_{\min}\right)v_j^p,
\qquad
p>1
\]
The exponent \(p\) controls the concentration near the shortest maturity. The map is strictly increasing, preserves both endpoints and leaves the QTT cardinality unchanged. Its derivative is
\[
\frac{\mathrm{d}T}{\mathrm{d}v}
=
p\left(T_{\max}-T_{\min}\right)v^{p-1}
\]
and therefore vanishes at \(v=0\). The smallest maturity cells are consequently placed near \(T_{\min}\), while their width increases progressively towards \(T_{\max}\).\\
\\
The first physical cell has width
\[
h_{T,\min}
=
\frac{T_{\max}-T_{\min}}
     {(n_T-1)^p}
\]
whereas the spacing of a uniform maturity grid is
\[
h_{T,\mathrm{unif}}
=
\frac{T_{\max}-T_{\min}}
     {n_T-1}
\]
Hence, the short-maturity refinement factor is
\[
\frac{h_{T,\min}}{h_{T,\mathrm{unif}}}
=
(n_T-1)^{1-p}
\]
In our experiments, we use the quadratic transformation
\[
p=2
\]
with \(n_T=8\). The first maturity cell is therefore seven times finer than under uniform discretization, without adding any QTT mode or increasing the number of tensor entries.\\
\\
This refinement complements the hyperbolic-sine moneyness grid. The latter concentrates nodes across the spatial direction normal to the payoff kink, whereas the maturity transformation resolves the rapid contraction of the convexity region as expiry approaches. The resulting two-dimensional node allocation is concentrated around the joint critical region
\[
m\simeq 0,
\qquad
T\simeq T_{\min}
\]
For American options, the same refinement also improves the representation of the short-dated exercise premium and of the rapidly evolving exercise frontier. It does not, however, eliminate the statistical noise of the LSMC pricing oracle. The benefit of the maturity grid must therefore be assessed separately from Monte Carlo label variance through the controlled ablation protocol introduced below.
\subsection{Preservation of the QTT architecture}
QTT tensorization acts on discrete grid indices rather than on the corresponding physical coordinates. For each dimension \(k\), the physical node is obtained through an index-to-coordinate map as describe in Eq.~\eqref{eq:grid_nodes}. 
The proposed transformations replace the affine maps of the moneyness and maturity axes by the nonlinear maps introduced in Sections 4.2 and 4.3. They preserve the number of nodes along every dimension, which remains a power of two. Consequently, the tensor shape, binary mode structure, mode ordering and admissible TT-cross index set are unchanged.\\
\\
When TT-cross selects a multi-index, the corresponding physical moneyness and maturity are first obtained from the transformed grids. The strike is then reconstructed as \(K=B(\boldsymbol{s})e^m\), and the pricing oracle returns
\begin{equation}
\mathcal{A}_{\boldsymbol{i}_S,j,q,\ell}
=
V
\left(
\boldsymbol{s}_{\boldsymbol{i}_S},
B(\boldsymbol{s}_{\boldsymbol{i}_S})e^{m_j},
r_q,
T_\ell
\right)
\end{equation}
TT-cross therefore requires no modification. Off-grid interpolation only uses the actual non-uniform node locations when computing the local weights, while the subsequent QTT contractions remain identical.\\
\\
Preserving the architecture does not imply preserving the ranks. The coordinate transformations change the numerical organization of the sampled price surface and may therefore change the QTT ranks required to reach a prescribed tolerance. Their effect on compressibility is assessed empirically through the rank, error and evaluation-budget comparisons reported in Section 6.
\section{Error Decomposition and Validation Protocol}
\label{sec:error-validation}
The deterministic decomposition
\eqref{eq:deterministic-error-decomposition} separates interpolation from
nodal reconstruction. We now extend it to a possibly biased and noisy pricing
oracle, and then specify how the resulting terms are isolated in the
experiments.
\subsection{Sources of Approximation Error}
\label{subsec:error-sources}
Let \(\mathcal{A}\) be the exact nodal pricing tensor defined in
Section~\ref{subsec:tt-qtt}, let \(\mathcal{Y}\) denote the tensor returned by
the pricing oracle, and let \(\widehat{\mathcal{Y}}\) be its TT-cross
reconstruction. Taking expectations entrywise over repeated oracle
realizations, define
\begin{equation}
\mathcal{B}_{\mathrm{orc}}
=
\mathcal{A}
-
\mathbb{E}\left[\mathcal{Y}\right]
\qquad
\mathcal{N}_{\mathrm{orc}}
=
\mathbb{E}\left[\mathcal{Y}\right]
-
\mathcal{Y}
\qquad
\mathbb{E}\left[\mathcal{N}_{\mathrm{orc}}\right]=0
\label{eq:oracle-bias-noise}
\end{equation}
where \(\mathcal{B}_{\mathrm{orc}}\) is the oracle bias measured as exact value
minus expected oracle output, and \(\mathcal{N}_{\mathrm{orc}}\) is its
centered sampling noise. The off-grid surrogate is
\(\widehat V=\mathcal{I}_h\widehat{\mathcal{Y}}\), so
\begin{equation}
\begin{aligned}
V(\boldsymbol{x})-\widehat V(\boldsymbol{x})
={}&
\underbrace{
V(\boldsymbol{x})
-
\mathcal{I}_h\mathcal{A}(\boldsymbol{x})
}_{\varepsilon_{\mathrm{grid}}(\boldsymbol{x})}
+
\underbrace{
\mathcal{I}_h\mathcal{B}_{\mathrm{orc}}(\boldsymbol{x})
}_{\varepsilon_{\mathrm{bias}}(\boldsymbol{x})}
\\
&+
\underbrace{
\mathcal{I}_h\mathcal{N}_{\mathrm{orc}}(\boldsymbol{x})
}_{\varepsilon_{\mathrm{noise}}(\boldsymbol{x})}
+
\underbrace{
\mathcal{I}_h
\left(
\mathcal{Y}-\widehat{\mathcal{Y}}
\right)(\boldsymbol{x})
}_{\varepsilon_{\mathrm{TT}}(\boldsymbol{x})}
\end{aligned}
\label{eq:complete-error-decomposition}
\end{equation}
The first term is the interpolation error imposed by the physical grid. Its
dependence on local spacing and curvature was established in
Section~\ref{subsec:off-grid-interpolation}; it is unaffected by increasing the
TT-cross budget. The adaptive grid is designed specifically to reduce this
term near the short-maturity convexity ridge without increasing the tensor
cardinality.\\
\\
For the analytical European oracle,
\(\mathcal{B}_{\mathrm{orc}}=\mathcal{N}_{\mathrm{orc}}=0\). For the American
LSMC oracle, the bias term includes time discretization, regression
approximation and the suboptimal stopping rule, whereas the noise term reflects
finite-path sampling \cite{longstaff2001leastSquares,glasserman2004montecarlo}.
Under standard regularity conditions, this sampling component retains the
usual Monte Carlo scale
\(\mathcal{O}(N_{\mathrm{path}}^{-1/2})\), while the regression and stopping
biases require separate convergence checks.
A fixed seed and common random numbers make \(\mathcal{Y}\) a consistent tensor
during one reconstruction, as discussed in
Section~\ref{subsec:reconstruction-conditions}, but they do not remove either
source of error.\\
\\
The final term measures how accurately TT-cross reproduces the particular
oracle tensor that it observes \cite{oseledets2010ttcross}. With noisy labels,
a small \(\varepsilon_{\mathrm{TT}}\) may therefore indicate that the
reconstruction fits a fixed noise realization rather than the exact price
surface. The four terms in \eqref{eq:complete-error-decomposition} are not
assumed to be independent, so their squared magnitudes cannot in general be
added. In practice, a budget-dependent error decrease identifies a
reconstruction-limited regime, whereas a persistent plateau indicates a grid
or oracle floor.
\subsection{Independent Validation and Metrics}
\label{subsec:validation-metrics}

All configurations are evaluated on the same set of off-grid states
\[
\mathcal{X}_{\mathrm{test}}
=
\left\{
\boldsymbol{x}^{(1)},\ldots,\boldsymbol{x}^{(N_{\mathrm{test}})}
\right\}
\]
which is generated independently of the TT-cross indices. Uniform and adaptive
grids use the same tensor shape, oracle budget, stopping tolerance and paired
reconstruction seeds. The GPR benchmark is evaluated on the same test set and
under the same oracle-evaluation budget.\\
\\
For European options, the reference labels are analytical and therefore isolate
the grid and reconstruction terms. For American options, they are obtained
from independent, higher-precision LSMC simulations using path sets and seeds
that were not used during reconstruction. One-dimensional binomial checks and
path and exercise-date convergence tests are used to verify that the remaining
reference error is small relative to the surrogate errors
\cite{longstaff2001leastSquares,glasserman2004montecarlo}.\\
\\
To compare states with different basket scales, define the normalized test
error
\begin{equation}
e_n
=
\frac{
\widehat V\left(\boldsymbol{x}^{(n)}\right)
-
V_{\mathrm{ref}}\left(\boldsymbol{x}^{(n)}\right)
}{
B\left(\boldsymbol{s}^{(n)}\right)
}
\label{eq:normalized-test-error}
\end{equation}
and report
\begin{equation}
\operatorname{MAE}
=
\frac{1}{N_{\mathrm{test}}}
\sum_{n=1}^{N_{\mathrm{test}}}|e_n|
\qquad
\operatorname{RMSE}
=
\left(
\frac{1}{N_{\mathrm{test}}}
\sum_{n=1}^{N_{\mathrm{test}}}e_n^2
\right)^{1/2}
\qquad
E_{\infty}
=
\max_n|e_n|
\label{eq:error-metrics}
\end{equation}
Percentage errors relative to the option price are not used as primary metrics
because they become unstable for prices close to zero. The same metrics are
also computed separately near the at-the-money region, at short maturities and
in the moneyness tails.\\
\\
Predictive accuracy is reported together with the effective QTT ranks, number
of oracle evaluations and wall-clock time. Comparisons across grid
configurations are paired by seed and repeated over increasing TT-cross
budgets. This protocol distinguishes a genuine reduction of the interpolation
floor from changes caused by initialization, computational effort or a
particular realization of the LSMC noise.

\section{Results}
\label{sec:results}
The experiments retain the eight-dimensional domain, tensor shape
\(32^5\times64\times8\times8\) and 37-core QTT representation used by
Gribben et al.\ \cite{Gribben2026STNGPR}. Four physical grids are considered:
\[
\begin{aligned}
\mathcal{G}_{00}&:\ \text{uniform absolute strike and uniform maturity}
\\
\mathcal{G}_{01}&:\ \text{uniform absolute strike and refined maturity}
\\
\mathcal{G}_{10}&:\ \text{hyperbolic-sine log-moneyness and uniform maturity}
\\
\mathcal{G}_{11}&:\ \text{hyperbolic-sine log-moneyness and refined maturity}
\end{aligned}
\]
They share the same tensor shape, QTT encoding and pricing domain; only the
mapping from discrete indices to physical strike or moneyness and maturity
values is modified.\\
\\
The experimental design is staged according to the purpose of each benchmark.
For the analytical European option, the original grid
\(\mathcal{G}_{00}\) is compared with the complete adaptive construction
\(\mathcal{G}_{11}\) over increasing TT-cross budgets, thereby measuring the
combined effect of the proposed coordinate transformations. For the American
option, a preliminary low-budget \(2\times2\) screening compares all four
configurations and identifies log-moneyness as the dominant source of
improvement. The main robustness experiment then compares
\(\mathcal{G}_{10}\) with \(\mathcal{G}_{11}\) across three budgets and five
paired reconstruction seeds, isolating the additional contribution of
short-maturity refinement after the moneyness transformation has been fixed.
All paired comparisons use identical test states and numerical settings; the
American errors are evaluated against independently generated,
higher-precision LSMC references.\\
\\
The comparison is therefore between different physical parameterizations of
the same discrete tensor architecture rather than tensors of different sizes.
The original study reported error and training-time curves for STN-GPR and GPR
on the uniform absolute-strike grid, but did not provide a coordinate ablation
or the associated effective ranks. The numerical baseline below is our
reproduction of that grid rather than a digitization of the published curves.
Unless stated otherwise, MAE values are expressed in option price units.
\subsection{European options}
\label{subsec:european-results}
The controlled benchmark is the European put on the geometric average of five
assets. All marginal volatilities are \(20\%\), all off-diagonal correlations
are \(0.3\), and the reference prices are analytical. The original and adaptive
grids are evaluated on the same \(2{,}000\) independently sampled off-grid
states. Since the tensor shape and TT-cross budgets are unchanged, the number
of oracle calls is nearly identical in every paired run.
\begin{table}[H]
    \centering
    \resizebox{\textwidth}{!}{%
    \begin{tabular}{rrrrrrrrr}
        \toprule
        & \multicolumn{2}{c}{Oracle evaluations}
        & \multicolumn{2}{c}{Effective rank}
        & \multicolumn{2}{c}{Price MAE}
        & \multicolumn{2}{c}{Prediction time (ms)} \\
        Budget & Original & Adaptive & Original & Adaptive
        & Original & Adaptive & Original & Adaptive \\
        \midrule
        \(5{,}000\)  & \(6{,}992\)  & \(6{,}994\)  & \(8.55\)  & \(4.47\) & \(4.622\) & \(0.104\)  & \(2.384\)  & \(0.949\) \\
        \(10{,}000\) & \(11{,}940\) & \(11{,}995\) & \(11.43\) & \(5.16\) & \(1.859\) & \(0.0845\) & \(3.729\)  & \(1.097\) \\
        \(20{,}000\) & \(21{,}942\) & \(21{,}803\) & \(15.29\) & \(6.08\) & \(1.426\) & \(0.0832\) & \(5.895\)  & \(1.446\) \\
        \(50{,}000\) & \(51{,}642\) & \(51{,}847\) & \(21.45\) & \(7.35\) & \(0.648\) & \(0.0831\) & \(10.747\) & \(1.799\) \\
        \bottomrule
    \end{tabular}}
    \caption{European geometric-basket put results. Prediction time is the
    measured mean wall-clock time per off-grid query. The TT-ANOVA
    initialization evaluations are included in the oracle counts.}
    \label{tab:european-grid-results}
\end{table}

\begin{figure}[H]
    \centering
    \includegraphics[width=\textwidth]{figures/results-european-grid.pdf}
    \caption{European effective rank and out-of-sample price MAE as functions
    of the nominal TT-cross budget. Both grids have the same tensor
    cardinality and are evaluated on the same analytical test set.}
    \label{fig:european-grid-results}
\end{figure}
The adaptive parameterization improves both accuracy and compressibility over
the complete budget range. At the largest budget, the effective rank decreases
from \(21.45\) to \(7.35\), a reduction of \(65.7\%\), while the price MAE
decreases from \(0.648\) to \(0.0831\). The corresponding normalized MAE
decreases from \(1.57\times10^{-2}\) to \(2.02\times10^{-3}\). The lower rank
also reduces the measured prediction time by \(83.3\%\), consistently with the
rank dependence of TT storage and contraction costs established in
Section~\ref{subsec:tt-qtt}.\\
\\
The gain is not merely a displacement along the same budget curve. With only
\(6{,}994\) oracle evaluations, the adaptive grid attains a MAE of \(0.104\),
whereas the original grid remains at \(0.648\) after \(51{,}642\) evaluations.
Within the tested range, the adaptive construction therefore achieves a lower
error with \(86.5\%\) fewer oracle calls than the most accurate reproduced
uniform-grid tensor. \\
\\
The error curve also identifies the limiting term in
\eqref{eq:complete-error-decomposition}. Exact multilinear interpolation of the
analytical nodal tensor on the adaptive grid has MAE \(0.08310\). The adaptive
TT reaches \(0.08323\) at budget \(20{,}000\) and \(0.08310\) at budget
\(50{,}000\). Further TT-cross sampling therefore removes essentially no test
error: reconstruction is no longer limiting and the remaining discrepancy is
the interpolation floor of the chosen physical grid. The original-grid curve
is still decreasing at the largest tested budget, so its interpolation floor
is not identified by these runs.\\
\\
These results extend, rather than contradict, the conclusion of the previous
study. Gribben et al.\ showed that TT-cross can access much larger effective
training sets than exact GPR on the original grid
\cite{Gribben2026STNGPR}. The present experiment shows that the physical
organization of those grid points is equally important: changing coordinates
without enlarging the tensor produces a smoother, substantially lower-rank
QTT representation.
\subsection{American options}
\label{subsec:american-results}
The American benchmark is a put on the arithmetic average of the same five
assets. Training labels use LSMC with \(2{,}000\) paths and 15 exercise dates.
Independent references use \(25{,}000\) paths, 60 exercise dates and five seeds
\(901,\ldots,905\). The robustness experiment contains \(200\) test states,
stratified across five moneyness and four maturity regimes, and compares
hyperbolic-sine log-moneyness with either uniform or short-maturity-refined
time nodes. Three budgets and five paired TT initializations give 30 completed
reconstructions. \\
\\
The pricing oracle was checked separately before comparing grids. Across 60
one-dimensional American-put scenarios, the \(25{,}000\)-path, 60-date LSMC
reference has mean absolute bias \(0.01365\) against a 3,000-step binomial tree,
with maximum absolute bias \(0.07449\). The binomial price lies inside the LSMC
\(95\%\) confidence interval in \(71.7\%\) of the cases, indicating that
regression and time-discretization bias remain visible. On the multi-asset
convergence set, the \(2{,}000\)-path, 15-date training oracle has MAE
\(0.01344\) and RMSE \(0.02721\) against the high-precision reference. Oracle
error is therefore smaller than the surrogate errors reported below, but it is
not negligible.
\begin{table}[H]
    \centering
    \resizebox{\textwidth}{!}{%
    \begin{tabular}{rrrrrrrr}
        \toprule
        & \multicolumn{2}{c}{Price MAE}
        & \multicolumn{2}{c}{Effective rank}
        & & \multicolumn{2}{c}{Fit time (s)} \\
        Budget & Uniform \(T\) & Adaptive \(T\) & Uniform \(T\) & Adaptive \(T\)
        & Adaptive wins & Uniform \(T\) & Adaptive \(T\) \\
        \midrule
        \(7{,}500\)  & \(0.1177\pm0.0647\) & \(0.0621\pm0.0119\) & \(10.20\pm0.34\) & \(10.07\pm0.18\) & \(5/5\) & \(68.0\) & \(58.7\) \\
        \(9{,}000\)  & \(0.1040\pm0.0675\) & \(0.0555\pm0.0093\) & \(11.07\pm0.46\) & \(10.95\pm0.19\) & \(5/5\) & \(78.2\) & \(69.2\) \\
        \(12{,}000\) & \(0.0992\pm0.0687\) & \(0.0582\pm0.0136\) & \(12.45\pm0.38\) & \(12.47\pm0.19\) & \(3/5\) & \(91.5\) & \(85.1\) \\
        \bottomrule
    \end{tabular}}
    \caption{American arithmetic-basket put robustness results. Entries are
    means and standard deviations across five paired TT initializations.
    Adaptive wins count the seeds for which short-maturity refinement lowers
    the MAE relative to uniform maturity at the same budget.}
    \label{tab:american-grid-results}
\end{table}

\begin{figure}[H]
    \centering
    \includegraphics[width=\textwidth]{figures/results-american-grid.pdf}
    \caption{American price MAE and effective rank across five paired
    reconstructions. Error bars show one standard deviation. Both curves use
    adaptive log-moneyness; only the maturity discretization differs.}
    \label{fig:american-grid-results}
\end{figure}
Short-maturity refinement lowers the mean MAE by \(47.2\%\), \(46.6\%\) and
\(41.3\%\) at budgets \(7{,}500\), \(9{,}000\) and \(12{,}000\), respectively.
It also reduces the dispersion across initializations by more than \(80\%\) at
each budget. The result is strongest at \(9{,}000\): all five paired runs
improve, the mean MAE is \(0.0555\), and the standard deviation is \(0.0093\). \\
\\
The rank evidence is more limited than in the European experiment. Mean ranks
are marginally lower at \(7{,}500\) and \(9{,}000\), and essentially unchanged
at \(12{,}000\). The lower mean fit times are consistent with a modest numerical
benefit, but they do not establish a structural complexity reduction of the
same magnitude as for the analytical European surface. The defensible
American conclusion is therefore improved accuracy and initialization
stability at comparable rank and oracle budget.\\
\\
Increasing the budget beyond \(9{,}000\) does not improve the adaptive mean
error. At \(12{,}000\), the rank and fit time increase, the MAE rises slightly,
and only three of the five paired seeds outperform uniform maturity. This point
fails the pre-specified \(80\%\) paired-win robustness criterion. The
\(9{,}000\)-evaluation regime is consequently the best observed compromise,
not a universally optimal budget. \\
\\
A preliminary full \(2\times2\) ablation at budget \(2{,}000\) gives MAE
\(9.817\), \(8.609\), \(0.903\) and \(0.394\) for
\(\mathcal{G}_{00},\mathcal{G}_{01},\mathcal{G}_{10}\) and
\(\mathcal{G}_{11}\), respectively. It indicates that log-moneyness supplies
the dominant improvement and that maturity refinement provides an additional
gain. Because this screening run uses only ten test points and a
\(500\)-path LSMC oracle, it is treated as qualitative support rather than as
the primary numerical result. \\
\\
The original STN-GPR paper used \(10{,}000\) LSMC paths and 30 exercise dates
and reported STN-versus-GPR accuracy as a function of training time
\cite{Gribben2026STNGPR}. Our multi-seed ablation deliberately uses a cheaper
training oracle and an independent, higher-precision reference. Absolute
wall-clock times and MAEs are therefore not claimed to be directly comparable
to its plotted values. The comparable object is the grid architecture: at the
same nominal tensor size, the transformed grid allocates the available TT-cross
samples more effectively in the regions that dominate American pricing error.

\section{Conclusion}
This work shows that the efficiency of a Tensor-Train pricing surrogate depends
not only on the reconstruction algorithm, but also on the physical coordinates
used to organize the pricing tensor. Building on the STN-GPR architecture of
Gribben et al.\ \cite{Gribben2026STNGPR}, we replaced absolute strike by basket
log-moneyness and concentrated a fixed number of nodes around the at-the-money
region and short maturities without modifying QTT architecture. For the
European benchmark, this reduces the effective rank from \(21.45\) to \(7.35\),
the price MAE from \(0.648\) to \(0.0831\), and the measured prediction time by
\(83.3\%\). The surrogate also reaches the interpolation floor of the adaptive
grid.\\
\\
For the American benchmark, short-maturity refinement reduces the mean error
and its dispersion across reconstruction seeds while leaving the effective
ranks broadly unchanged. Budget \(9{,}000\) provides the best observed
compromise; a larger budget increases cost without improving mean accuracy.
The American evidence therefore supports improved accuracy and robustness, but
not a general claim of rank reduction, and the attainable error remains partly
limited by the bias and sampling noise of the LSMC oracle.\\
\\
These results establish coordinate design as a structural component of
tensor-network option pricing. A second paper will use the framework developed
here to investigate Delta, Gamma and cross-Gamma estimation, and to assess the
surrogate as a repeated revaluation engine for Value-at-Risk and Expected
Shortfall. Particular attention will be paid to error amplification under
differentiation and to the validation of portfolio-tail estimates. Separating
these questions keeps the present contribution focused on the pricing,
compression and robustness properties on which those applications depend.\\
\\
The accuracy gained through this coordinate design, particularly for European baskets, makes the resulting grid possibly usable for fast pricing in a trading context. A natural use case is the computation of Greeks making the surrogate well suited to intraday hedging and to model validation workflows where dense and repeated revaluation is required.\\
\\
Beyond its immediate classical benefits, the proposed QTT representation also provides a natural interface with future quantum implementations. Tensor Trains are mathematically equivalent to Matrix Product States, a standard tensor-network representation of multiqubit quantum states, and low-rank MPS representations can be compiled into quantum circuits composed of one- and two-qubit gates \cite{orus2014tensorNetworks,Ran2020MPSQuantumCircuits}. By expressing the pricing surface through binary modes while preserving low tensor ranks, the present framework therefore produces a representation that could serve as a structured starting point for quantum state preparation or hybrid quantum--classical pricing algorithms. This connection does not imply an automatic quantum advantage, since normalization, data loading, circuit depth, measurement costs and hardware noise would still need to be addressed. Nevertheless, it makes the proposed implementation compatible with a longer-term transition towards quantum computing, without requiring the high-dimensional pricing problem to be reformulated from scratch.

\appendix

\section{Proof of Separability}
\label{app:separability-proof}
Let $\Omega$ be the computational space, formally
\[
\Omega = I_{S_1}\times \dots \times I_{S_5} \times I_m \times I_r \times I_T \subset \mathbb{R}^8
\]
where 
\[
I_S = [S_{min},S_{max}], \quad I_m = [m_{min}, m_{max}], \quad I_r = [r_{min}, r_{max}], \quad I_T =[T_{min}, T_{max}]
\]
or more generally 
\[
\Omega = \prod _{k=1} ^d I_k, \qquad d = d_S +3
\]
For each interval, $I_k = [a_k, b_k]$, consider 
\[
D_k = \left(\mathbb{Q} \cap (a_k,b_k) \right) \cup \{a_k,b_k\}
\]
The set $D_k$ is countable because $\mathbb{Q}$ is countable. $\D_k$ is also dense in $I_k$. Let define 
\[
D = D_1 \times \dots \times D_d \subset \Omega
\]
As $D$ is a finite product of countable set it is countable. Now let's show that $D$ is dense in $\Omega$. Let 
\[
\boldsymbol{x} = (x_1,\dots,x_d) \in \Omega
\]
and let's $\varepsilon > 0$. Since $D_k$ is dense in $I_k$, there exist for every $k$, an element $q_k \in D_k$ such that :
\[
|x_k - q_k| < \frac{\varepsilon}{\sqrt{d}}
\]
defining $\boldsymbol{q} = (q_1,\dots,q_k) \in D$, we obtain 
\[
||\boldsymbol{x} - \boldsymbol{q}||_2 ^2 = \sum _{k=1} ^d |x_k - q_k|^2 < \sum ^d _{k=1} \frac{\varepsilon ^2}{d} = \varepsilon^2
\]
So finally 
\[
||\boldsymbol{x} - \boldsymbol{q}||_2 < \varepsilon
\]
So any open ball centered on a point of $\Omega$ contain an element of $D$. Consequently 
\[
\bar{D} = D
\]
$\Omega$ possess a dense countable part, so $\Omega$ is separable by definition. It appear, thank to this demonstration that adding $I_{\sigma}$ does not annihilate the separability of the space in a topological sense. 

\newpage
\bibliographystyle{plain}
\bibliography{references}

\end{document}



